\documentclass[11pt]{article}
\usepackage[margin=1.1in]{geometry}
\usepackage{amsmath,amssymb,amsthm,mathtools,booktabs,array,tabularx,longtable}
\usepackage[dvipsnames]{xcolor}
\usepackage[colorlinks=true,linkcolor=MidnightBlue,citecolor=MidnightBlue,urlcolor=MidnightBlue]{hyperref}
\usepackage{fancyhdr}
\pagestyle{fancy}\fancyhf{}\renewcommand{\headrulewidth}{0pt}
\fancyfoot[R]{\thepage}
\fancyfoot[L]{\scriptsize\color{gray}Generated by the Transport Lab from approved parts of its store; not a deposited paper. Papers of record: \href{https://zenodo.org/communities/transport-papers}{Zenodo, Transport Papers}}
\fancypagestyle{plain}{\fancyhf{}\fancyfoot[R]{\thepage}\fancyfoot[L]{\scriptsize\color{gray}Generated by the Transport Lab from approved parts of its store; not a deposited paper. Papers of record: \href{https://zenodo.org/communities/transport-papers}{Zenodo, Transport Papers}}\renewcommand{\headrulewidth}{0pt}}
\newtheorem{theorem}{Theorem}[section]
\newtheorem{proposition}[theorem]{Proposition}
\newtheorem{lemma}[theorem]{Lemma}
\newtheorem{corollary}[theorem]{Corollary}
\newtheorem{conjecture}[theorem]{Conjecture}
\newtheorem{refutation}[theorem]{Refutation}
\newtheorem{prediction}[theorem]{Prediction}
\newtheorem{principle}[theorem]{Principle}
\newtheorem{observation}[theorem]{Observation}
\newtheorem{axiom}[theorem]{Axiom}
\theoremstyle{definition}
\newtheorem{definition}[theorem]{Definition}
\newtheorem{fact}[theorem]{Fact}
\newtheorem{assembled}[theorem]{Assembly}
\newtheorem{postulate}[theorem]{Postulate}
\newtheorem{premise}[theorem]{Premise}
\newtheorem{hypothesis}[theorem]{Hypothesis}
\newtheorem{remark}[theorem]{Remark}
\newtheorem{example}[theorem]{Example}
\theoremstyle{plain}
\setlength{\marginparwidth}{0.8in}
\newcommand{\changedmark}[2]{\marginpar{\raggedright\scriptsize\color{#1}#2}}
\providecommand{\Aa}{}\renewcommand{\Aa}{\mathcal{A}}
\providecommand{\Ad}{}\renewcommand{\Ad}{\widetilde{\mathrm{Ad}}}
\providecommand{\Aev}{}\renewcommand{\Aev}{\mathcal{A}^{+}}
\providecommand{\Aod}{}\renewcommand{\Aod}{\mathcal{A}^{-}}
\providecommand{\B}{}\renewcommand{\B}{\mathrm{B}}
\providecommand{\C}{}\renewcommand{\C}{\mathbb{C}}
\providecommand{\CAR}{}\renewcommand{\CAR}{\mathrm{CAR}}
\providecommand{\CS}{}\renewcommand{\CS}{{\mathbb{S}^{+}}}
\providecommand{\CSm}{}\renewcommand{\CSm}{{\mathbb{S}^{-}}}
\providecommand{\Cl}{}\renewcommand{\Cl}{\mathrm{Cl}}
\providecommand{\Con}{}\renewcommand{\Con}{\mathcal{C}}
\providecommand{\Dl}{}\renewcommand{\Dl}{\Delta}
\providecommand{\EN}{}\renewcommand{\EN}{E_{N}}
\providecommand{\End}{}\renewcommand{\End}{\mathrm{End}}
\providecommand{\F}{}\renewcommand{\F}{\mathbb{F}}
\providecommand{\G}{}\renewcommand{\G}{\Gamma}
\providecommand{\HH}{}\renewcommand{\HH}{\mathbb{H}}
\providecommand{\IS}{}\renewcommand{\IS}{I\sigma_3}
\providecommand{\J}{}\renewcommand{\J}{\mathbf{J}}
\providecommand{\LH}{}\renewcommand{\LH}{L_{H}}
\providecommand{\Mat}{}\renewcommand{\Mat}{\mathrm{M}}
\providecommand{\Mo}{}\renewcommand{\Mo}{M^{\circ}}
\providecommand{\Mp}{}\renewcommand{\Mp}{M^{+}}
\providecommand{\Msun}{}\renewcommand{\Msun}{M_{\odot}}
\providecommand{\Mt}{}\renewcommand{\Mt}{\widetilde{M}}
\providecommand{\Ncone}{}\renewcommand{\Ncone}{\mathcal{N}}
\providecommand{\Neff}{}\renewcommand{\Neff}{N_{\mathrm{eff}}}
\providecommand{\PCS}{}\renewcommand{\PCS}{\mathbb{P}{\mathbb{S}^{+}}}
\providecommand{\PT}{}\renewcommand{\PT}{\mathbb{PT}}
\providecommand{\Pin}{}\renewcommand{\Pin}{\mathrm{Pin}}
\providecommand{\Q}{}\renewcommand{\Q}{\mathbb{Q}}
\providecommand{\R}{}\renewcommand{\R}{\mathbb{R}}
\providecommand{\Ric}{}\renewcommand{\Ric}{\mathrm{Ric}}
\providecommand{\SBH}{}\renewcommand{\SBH}{S_{\mathrm{BH}}}
\providecommand{\Sc}{}\renewcommand{\Sc}{\mathbb{S}}
\providecommand{\Sch}{}\renewcommand{\Sch}{\mathrm{P}}
\providecommand{\Sig}{}\renewcommand{\Sig}{\Sigma}
\providecommand{\Som}{}\renewcommand{\Som}{S_\omega}
\providecommand{\Sp}{}\renewcommand{\Sp}{\mathbb{S}}
\providecommand{\Sph}{}\renewcommand{\Sph}{\mathrm{S}}
\providecommand{\Spi}{}\renewcommand{\Spi}{S_\pi}
\providecommand{\Spin}{}\renewcommand{\Spin}{\mathrm{Spin}}
\providecommand{\Sth}{}\renewcommand{\Sth}{S_{\mathrm{th}}}
\providecommand{\T}{}\renewcommand{\T}{\mathbb{T}}
\providecommand{\TH}{}\renewcommand{\TH}{T_{\mathrm{H}}}
\providecommand{\Tr}{}\renewcommand{\Tr}{\operatorname{Tr}}
\providecommand{\UHF}{}\renewcommand{\UHF}{\mathrm{UHF}}
\providecommand{\Z}{}\renewcommand{\Z}{\mathbb{Z}}
\providecommand{\Zc}{}\renewcommand{\Zc}{\mathrm{Z}}
\providecommand{\cmg}{}\renewcommand{\cmg}{\mathrm{cm^{2}\,g^{-1}}}
\providecommand{\dS}{}\renewcommand{\dS}{\mathrm{dS}}
\providecommand{\dsc}{}\renewcommand{\dsc}{\mathrm{dsc}}
\providecommand{\dtPG}{}\renewcommand{\dtPG}{dt}
\providecommand{\dtS}{}\renewcommand{\dtS}{dt_{\mathrm{S}}}
\providecommand{\ellv}{}\renewcommand{\ellv}{\ell}
\providecommand{\er}{}\renewcommand{\er}{e_r}
\providecommand{\et}{}\renewcommand{\et}{e_t}
\providecommand{\ev}{}\renewcommand{\ev}{\mathrm{ev}}
\providecommand{\gl}{}\renewcommand{\gl}{\mathfrak{gl}}
\providecommand{\hatV}{}\renewcommand{\hatV}{\widehat{V}}
\providecommand{\hb}{}\renewcommand{\hb}{\bar{h}}
\providecommand{\hf}{}\renewcommand{\hf}{\hat{f}}
\providecommand{\id}{}\renewcommand{\id}{\mathrm{id}}
\providecommand{\kpc}{}\renewcommand{\kpc}{\mathrm{kpc}}
\providecommand{\lP}{}\renewcommand{\lP}{\ell_{P}}
\providecommand{\mD}{}\renewcommand{\mD}{\mathbf{m}_{D}}
\providecommand{\mb}{}\renewcommand{\mb}{\bar m}
\providecommand{\mnu}{}\renewcommand{\mnu}{\Sigma m_{\nu}}
\providecommand{\nInf}{}\renewcommand{\nInf}{n_\infty}
\providecommand{\nO}{}\renewcommand{\nO}{n_0}
\providecommand{\ninf}{}\renewcommand{\ninf}{n_{\infty}}
\providecommand{\nz}{}\renewcommand{\nz}{n_{0}}
\providecommand{\proj}{}\renewcommand{\proj}{\operatorname{proj}}
\providecommand{\pv}{}\renewcommand{\pv}{\pi}
\providecommand{\rad}{}\renewcommand{\rad}{\mathrm{rad}}
\providecommand{\sgn}{}\renewcommand{\sgn}{\operatorname{sgn}}
\providecommand{\shA}{}\renewcommand{\shA}{\mathcal{A}}
\providecommand{\shC}{}\renewcommand{\shC}{\mathcal{C}}
\providecommand{\shE}{}\renewcommand{\shE}{\mathcal{E}}
\providecommand{\shJ}{}\renewcommand{\shJ}{\mathcal{J}}
\providecommand{\shK}{}\renewcommand{\shK}{\mathcal{K}}
\providecommand{\shL}{}\renewcommand{\shL}{L}
\providecommand{\shW}{}\renewcommand{\shW}{\mathcal{W}}
\providecommand{\sll}{}\renewcommand{\sll}{\mathfrak{sl}}
\providecommand{\so}{}\renewcommand{\so}{\mathfrak{so}}
\providecommand{\su}{}\renewcommand{\su}{\mathfrak{su}}
\providecommand{\supp}{}\renewcommand{\supp}{\operatorname{supp}}
\providecommand{\tr}{}\renewcommand{\tr}{\operatorname{tr}}
\providecommand{\uu}{}\renewcommand{\uu}{\mathfrak{u}}
\providecommand{\br}{}\renewcommand{\br}[1]{\langle #1\rangle}
\providecommand{\bra}{}\renewcommand{\bra}[1]{\langle #1|}
\providecommand{\file}{}\renewcommand{\file}[1]{\texttt{#1}}
\providecommand{\hp}{}\renewcommand{\hp}[2]{( #1 \,|\, #2 )}
\providecommand{\ip}{}\renewcommand{\ip}[2]{\langle #1, #2\rangle}
\providecommand{\ket}{}\renewcommand{\ket}[1]{|#1\rangle}
\providecommand{\rel}{}\renewcommand{\rel}[2]{D\bigl(#1 \,\|\, #2\bigr)}
\providecommand{\rev}{}\renewcommand{\rev}[1]{\widetilde{#1}}
\providecommand{\slashed}{}\renewcommand{\slashed}[1]{{#1}\mkern-10mu/}
\title{\bfseries Temperature and Hawking radiation}
\author{Leon Butler\\ \normalsize\href{mailto:lbutler.math@gmail.com}{\texttt{lbutler.math@gmail.com}}\\ \normalsize ORCID: \href{https://orcid.org/0009-0006-4449-5617}{0009-0006-4449-5617}}
\date{Generated by the Transport Lab, 2 November 2026}
\begin{document}
\maketitle
\begin{abstract}
\normalsize
A black hole glows at a temperature set by its mass, and this trail shows where that number comes from in the framework and what rests on it. The horizon's surface gravity, $\kappa = 1/4M$, is computed from the gravitational field.

Under premises (P) and (R), every boundary mode is thermal at the Hawking temperature. Assuming the interface premise (I), that thermal state reaches a distant observer as Hawking's spectrum with greybody factors.

The second half asks which state the horizon modes are in. The premise (N$_{+}$) says only that they are in a ground state for sliding along the light rays. Its short-distance form now holds for every Hadamard state of the Dirac field, by the theorem (HS). That the series' states are Hadamard is itself an adopted premise, (HD), which does not say which state it is.
\end{abstract}
\newpage
\noindent \textbf{Authorship and disclosure. }The founding idea of the series is the author's: that division by zero in a Clifford algebra should be read not as a prohibition but as a transport into a higher-dimensional algebra, and that the transported division need not be accompanied by an inverse operation. The direction of investigation and the critical review of the results are likewise the author's. The mathematical development and the drafting of the text were carried out in extended collaboration with Claude Fable 5 (Anthropic), an AI system.

\noindent\fbox{\parbox{\dimexpr\linewidth-2\fboxsep-2\fboxrule}{\textbf{This is not a deposited paper.} It was generated by the Transport Lab on 2 November 2026 and assembled from approved parts of its store. It follows results across several papers of the series; each is stated with its proof and credited to the paper that first proved it. The official versions of those papers are on Zenodo, in the Transport Papers community, \url{https://zenodo.org/communities/transport-papers}.}}
\newpage

\section*{Introduction}
The walk starts with the gravitational field, moves to the temperature of a single mode, and follows the radiation out to a distant observer. It then asks how much of the state behind that temperature follows from the field theory, and what is still assumed. Each result appears after the ones it rests on.

\section{What we assume}

\begin{postulate}[Premise (I): the interface identification]\label{c:hawking.premise-i-interface-identification}
This premise assumes that whatever the horizon state puts into each outgoing channel is carried out to distant observers in the ordinary way. It is an assumption, not a theorem of the framework.
\end{postulate}

\begin{postulate}[(P): univalence]\label{c:singularity.p-univalence}
This is an assumption, not a theorem. It says that the physical states of the modes living on the boundary give zero average to every quantity that changes fermion number by an odd amount. So expectation values never mix the even and odd parity sectors.
\end{postulate}

\begin{postulate}[(R): regularity at the fixed locus]\label{c:singularity.r-regularity-fixed-locus}
This is an assumption, not a theorem. It asks that each boundary mode, taken on its own, looks thermal with respect to the black hole's time flow. The temperature is fixed by the gravitational field: it is exactly the one for which the imaginary-time version of the field describes a smooth geometry, with no cone-shaped point.
\end{postulate}

\begin{postulate}[(HD), the Hadamard deposit premise: the descent's deposited state is a Hadamard state on a Cauchy collar of the descent, for the Dirac field and for the scalar fields of the descended phase; the occupation profile over an admissible reference vacuum stays free data]\label{c:synthesis.hadamard-deposit-premise}
Quantum fields in a gravitational field make sense only in well-behaved states, called Hadamard states, which look like empty space at very short distances. The series' account of the universe's earliest moment, the descent, leaves its fields in some state, but nothing in the series yet shows that state is well behaved. This premise assumes it: just after the descent the state is of the well-behaved kind. Since good behaviour, once present, is kept by the equations of motion, every later situation, including black holes formed by collapse, then has well-behaved states, which is what the series' horizon results and the energy of quantum fields need. The premise does not say which well-behaved state it is, how many particles of each energy were left behind, or what pattern of primordial ripples they carry; those remain open inputs.
\end{postulate}

\begin{postulate}[(N$_{+}$): on the continuum of a horizon's null generators the state of the modes is a gauge-invariant quasi-free ground state for the affine null translations along the generators]\label{c:synthesis.n-plus-null-ground-state}
Black holes, the horizon of an expanding universe and the horizon an accelerating observer sees all have a temperature. The series used to assume each temperature with its own premise. This premise replaces them with one assumption that mentions no temperature: the quantum modes running along each horizon's light rays are in their lowest-energy state for sliding along the rays. The store proves that this state, seen from one side of the horizon, is exactly thermal at the expected temperature. The assumption describes a black hole that has always existed in equilibrium; read on the light rays falling into the horizon, its short-distance form also covers a hole formed by collapse, given a simplification of the radiating fields to one direction per light ray. For the horizon an accelerating observer sees, and for the cosmological horizon with fields that keep their form under rescalings of lengths, it amounts to choosing the standard vacuum of the field theory. It still carries an input: that the state is smooth where the two halves of the horizon meet, the condition every derivation of black-hole temperature uses in some form. A black hole that has always existed in equilibrium is an idealisation in a universe with a positive cosmological constant, where no such equilibrium state exists, so the local form of the assumption is the one that matters there; whether the series' states are smooth enough at short distances to give that local form everywhere is still open, though it now comes down to the usual short-distance requirement on states: how such states look at very short distances where gravity is present, assumed before, is now proved.
\end{postulate}

\section{Temperature and Hawking radiation}

This section goes from the shape of the horizon to the light it emits, and then to the state of the modes that emit it.

Everything starts with a single number describing how strongly the horizon pulls. Surface gravity is the redshift-corrected acceleration needed to hover just outside the horizon; with $f = 1 - 2M/r$ the norm of the static Killing field, \[ \kappa = \lim\_{r \to 2M} \frac{f'(r)}{2} = \frac{1}{4M} \] in natural units.

This rate will set the temperature.

\begin{proposition}[$\kappa = 1/4M$, computed]\label{c:singularity.computed}
Surface gravity measures how strongly a black hole's horizon pulls, as judged from far away. It is found by taking the acceleration an observer needs to hover at fixed distance, correcting it for the gravitational redshift, and following that quantity down to the horizon. For a non-rotating black hole of mass M the result is 1/(4M) in natural units.
\emph{First proved in} T5, Prop A.2 (\href{https://doi.org/10.5281/zenodo.22721790}{doi:10.5281/zenodo.22721790}).
\end{proposition}
\noindent\emph{Rests on:} \ref{c:singularity.imported-surface-gravity-redshifted-acceleration}.

A rate becomes a temperature through a period, and the field supplies one. Replacing time by an imaginary variable $\tau$ closes the time direction into a circle, and near the horizon the section looks like a plane in polar coordinates with angle $\kappa\tau$.

It is smooth there only if the circle has length $2\pi/\kappa$; any other length leaves a cone point, so the period is fixed rather than chosen.

\begin{proposition}[The Euclidean section, and its period]\label{c:singularity.euclidean-section-period}
Replacing time by an imaginary variable turns the gravitational field outside a black hole into a field on a flat four-dimensional space without time, in which the time direction closes into a circle. Near the horizon the geometry this field describes looks like a flat plane in polar coordinates, but only if the circle has exactly the right length, $8\pi$ times the mass in natural units; any other length leaves a cone point. So the period is fixed by the gravitational field rather than chosen.
\emph{First proved in} T5, Prop 7.18 (\href{https://doi.org/10.5281/zenodo.22721790}{doi:10.5281/zenodo.22721790}).
\end{proposition}
\noindent\emph{Rests on:} \ref{c:singularity.computed}.

The next step connects this geometry to the algebra of the boundary modes. The curvature of the gravitational gauge field picks out, with no arbitrary choice, the plane of the boost bivector $\sigma_r$, and at the horizon the rate of that boost is the surface gravity $\kappa = 1/4M$.

That rate converts the algebra's natural angle into ordinary time: an interval of $2\pi$ in rapidity is $2\pi/\kappa = 8\pi M$ of Schwarzschild time. This single conversion factor is where the numerical value of the temperature enters.

\begin{proposition}[The plane and the rate]\label{c:singularity.plane-rate}
The curvature of a non-rotating black hole picks out, at every distance from the centre, a particular pair of planes without any arbitrary choice. At the horizon one of these is the plane in which the horizon's own time symmetry acts as a boost, and the rate of that boost is the surface gravity, one over four times the mass in natural units. That rate is the single number converting the geometry's natural angle into ordinary time, and it is where the numerical value of the black hole's temperature comes from.
\emph{First proved in} T5, Prop 6.2 (\href{https://doi.org/10.5281/zenodo.22721790}{doi:10.5281/zenodo.22721790}).
\end{proposition}
\noindent\emph{Rests on:} \ref{c:singularity.computed}, \ref{c:singularity.curvature-gauge-field-exact-shape-decoupling}.

One might hope the algebra alone could fix the temperature, and this result shows it cannot. A symmetry $\theta$ of a single mode's algebra that respects the flow fixes every equilibrium state, and one that reverses the flow only flips the sign: \[ \omega\_\beta \circ \theta = \omega\_{-\beta} \]

So the size $|\beta|$ of the inverse temperature depends on how the flow's parameter is converted into physical time, which is exactly what the geometric rate supplies.

\begin{theorem}[Rate rigidity: the temperature is not algebra intrinsic]\label{c:singularity.rate-rigidity-temperature-not-algebra-intrinsic}
For a single quantum mode with a given time evolution, the symmetries of the algebra that respect the time evolution, or reverse it, can at most flip the sign of an equilibrium state's inverse temperature; they never change its size. Rescaling the clock rescales the temperature while leaving the state itself alone. So the algebra on its own cannot pick out a particular temperature: that needs outside information, namely how the evolution's parameter is converted into physical time.
\emph{First proved in} T5, Thm 7.16 (\href{https://doi.org/10.5281/zenodo.22721790}{doi:10.5281/zenodo.22721790}).
\end{theorem}
\noindent\emph{Rests on:} \ref{c:singularity.kms-uniqueness-factor}.

With the rate in hand, the two premises in what we assume can be put to work. Under univalence (P) and regularity (R), every boundary mode is thermal at the same inverse temperature, \[ \beta = \frac{2\pi}{\kappa} = 8\pi M , \qquad T = \frac{1}{8\pi M} = \TH , \] with no assumption about correlations between modes.

As a by-product, the condition that each mode is entangled with its partner now follows from the premises instead of being assumed.

\begin{theorem}[The rate, and the discharge of hypothesis (iii)]\label{c:singularity.rate-discharge-hypothesis-iii}
Two physical premises, one restricting which states are allowed and one requiring smoothness of the geometry at the horizon, force each boundary mode of the black hole into a thermal state. Every mode comes out at the same temperature, the Hawking temperature, without any assumption about correlations between the modes. As a by-product, a condition that previously had to be assumed, that each mode is genuinely entangled with its partner, now follows from the premises.
\emph{First proved in} T5, Thm 7.20 (\href{https://doi.org/10.5281/zenodo.22721790}{doi:10.5281/zenodo.22721790}).
\end{theorem}
\noindent\emph{Rests on:} \ref{c:singularity.ascended-marginals-are-diagonal}, \ref{c:singularity.computed}, \ref{c:singularity.descended-form-state-independent}, \ref{c:singularity.modular-structure-completed-pair}, \ref{c:singularity.plane-rate}, \ref{c:singularity.p-univalence}, \ref{c:singularity.r-regularity-fixed-locus}.

Single modes are thermal, and the next result assembles them into one state for the whole boundary. The I-state fills every boundary mode $j$ of frequency $\omega_j$ by the Fermi-Dirac rule, \[ n\_j = \frac{1}{e\^{}{\omega\_j/\TH} + 1} \] and it is thermal in the precise algebraic sense, called KMS, on finite windows of modes and on the whole system.

Among the simple Gaussian, or quasi-free, states it is the only one thermal for the boost clock, so its occupations are derived rather than chosen.

\begin{theorem}[The I-state]\label{c:hawking.i-state}
The boundary modes in this black-hole setting can be placed in a thermal state at the Hawking temperature, each mode filled according to the Fermi-Dirac rule. This result shows that this state exists and is thermal in the precise algebraic sense, both on every finite set of modes and on the whole system. It also shows that among the simple Gaussian (quasi-free) states it is the only one that is thermal for the geometry's boost time evolution, which turns the phase of each mode at its own frequency, so its occupation numbers are derived rather than chosen.
\emph{First proved in} T7, Thm 3.6 (\href{https://doi.org/10.5281/zenodo.22722121}{doi:10.5281/zenodo.22722121}).
\end{theorem}
\noindent\emph{Rests on:} \ref{c:hawking.gibbs-states-are-kms-conversely-finite}, \ref{c:hawking.imported-kms-passes-inductive-limits}, \ref{c:hawking.imported-quasi-free-states-car-algebra}.

Real radiation comes in a continuous range of frequencies, not a countable list. On the continuum of modes, with one-particle generator $h$, the quasi-free state with symbol $A_I = (1 + e^{\beta h})^{-1}$ is KMS at $\beta = 2\pi/\kappa$.

It is the only KMS state, and no requirement that the energies be positive is needed.

\begin{theorem}[The I-state on the continuum]\label{c:hawking.i-state-continuum}
The earlier result about the horizon's equilibrium state was proved for a countable list of modes. This result carries it over to a continuous range of frequencies, as real radiation needs. The same thermal state, at the Hawking temperature, is again the only equilibrium state, and the energies need not be positive for this to hold.
\emph{First proved in} T7, Thm 4.2 (\href{https://doi.org/10.5281/zenodo.22722121}{doi:10.5281/zenodo.22722121}).
\end{theorem}
\noindent\emph{Rests on:} \ref{c:hawking.i-state}, \ref{c:hawking.imported-quasi-free-kms-reduces-two}, \ref{c:hawking.imported-quasi-free-states-car-algebra}, \ref{c:hawking.imported-ideal-fermi-gas-uniqueness}.

Black holes also emit bosons such as photons, so the fermionic picture must produce them. Under the s-wave modelling convention (Premise (Q)), with level spacing $\epsilon$, the collective modes $a_m$ built from the thermally filled fermion ladder satisfy, on the infinite ladder, \[ \sum\_{k \in \Z} n\_{k+m}\,(1 - n\_k) = \frac{m}{e\^{}{\beta\epsilon m} - 1} \] which is $m$ times the Bose-Einstein occupancy.

On a finite window the occupancy falls short by a positive amount that shrinks exponentially as the edges move away from the filled sea.

\begin{proposition}[Bose-Einstein occupancy, with quantified window corrections]\label{c:hawking.bose-einstein-occupancy-quantified-window-corrections}
A ladder of fermion levels, filled thermally in the Fermi-Dirac pattern, can be repackaged into collective boson modes. Given the s-wave modelling convention for the ladder, this result shows by an exact algebraic identity that on the infinite ladder those collective modes have exactly the Bose-Einstein thermal occupancy at the Hawking temperature. On a finite window of levels centred on the filled sea, the answer falls short only by a positive amount that shrinks exponentially as the window edges move away from the filled sea.
\emph{First proved in} T7, Prop 6.2 (\href{https://doi.org/10.5281/zenodo.22722121}{doi:10.5281/zenodo.22722121}).
\end{proposition}
\noindent\emph{Rests on:} \ref{c:hawking.i-state}, \ref{c:hawking.imported-s-wave-bosonic-sector}.

The last step carries the radiation from the horizon to a distant observer. Assuming the interface premise (I), that emission near the horizon is carried outward in the ordinary way, the flux in each channel is \[ \frac{d$\mathcal{N}$}{dt\,d\omega} = \frac{1}{2\pi}\,\Gamma\_{s\ell m}(\omega)\; n\_\pm(\omega) , \qquad n\_\pm(\omega) = \frac{1}{e\^{}{\omega/\TH} \pm 1} \] with $\Gamma_{s\ell m}$ the classical transmission through the potential barrier.

That is Hawking's spectrum with its usual greybody factors, the bosonic part resting also on premise (Q).

\begin{theorem}[Greybody inheritance, under the interface premise]\label{c:hawking.greybody-inheritance-under-interface-premise}
Radiation leaving a black hole is thermal at the horizon, but part of it is scattered back by the black hole's gravitational field outside the horizon, so what reaches a distant observer is filtered. This result combines the framework's thermal horizon state with the ordinary general-relativistic filtering, under the interface premise (I), that what is emitted near the horizon is carried outward in the standard way. For the boson-like part it also relies on a further premise, (Q), about the joint state of the modes. The outcome is the usual Hawking spectrum, including its filtering factors.
\emph{First proved in} T7, Asm 5.4 (\href{https://doi.org/10.5281/zenodo.22722121}{doi:10.5281/zenodo.22722121}).
\end{theorem}
\noindent\emph{Rests on:} \ref{c:hawking.imported-exterior-exactly-general-relativity}, \ref{c:hawking.i-state-continuum}, \ref{c:hawking.imported-s-wave-bosonic-sector}, \ref{c:hawking.imported-modular-descent-premises-geometric-normalisation}, \ref{c:hawking.premise-i-interface-identification}.

So far the temperature has come from premises that name it. The premise (N$_{+}$) replaces their thermal content with one statement that has no temperature in it: on the continuum of a horizon's null generators the modes are in a gauge-invariant quasi-free ground state for the affine null translations along the generators.

Seen from either side of the edge, such a state is exactly thermal at rapidity period $2\pi$, so the continuum forms of (R) and its companions follow from it. The thermal input is moved, not removed: it now sits in the regularity of the state across the bifurcation surface, where the two halves of the horizon meet.

The eternal hole is an idealisation when $\Lambda > 0$, so the premise's local form, near each point of a horizon, is the one that carries the physics. The rest of the trail asks how much of this premise the field theory supplies.

Physicists already have a standard short-distance requirement on quantum states, and the premise will be measured against it. A quasi-free state is Hadamard when the singular part of its two-point function carries only positive frequency, the wave front set condition $\mathrm{WF}'(\Lambda) = C^{+}$, imposed for the Dirac field on both orderings.

Any two Hadamard states differ by a smooth kernel. This is imported from the literature (Radzikowski for scalar fields, Hollands for the Dirac field), not proved in the series.

\begin{fact}[Imported: the microlocal form of the Hadamard condition (Radzikowski 1996), and for the Dirac field on both orderings of the two-point function (Hollands 2001), under which Hadamard states differ by smooth kernels]\label{c:synthesis.imported-microlocal-hadamard-condition}
Physicists require that a reasonable quantum state of a field, looked at closely enough, resembles empty space. Radzikowski found a precise way to say this in terms of where the state's correlations are infinitely sharp and in which direction their energy flows, and Hollands extended it to electrons and other particles of spin one half. Any two states meeting the requirement differ only by something smooth. These are taken from the literature.
\emph{Source:} M. J. Radzikowski, Micro-local approach to the Hadamard condition in quantum field theory on curved space-time, Commun. Math. Phys. 179 (1996), 529--553.; S. Hollands, The Hadamard condition for Dirac fields and adiabatic states on Robertson-Walker spacetimes, Commun. Math. Phys. 216 (2001), 635--661; arXiv:gr-qc/9906076., Def. IV.1; Thm. IV.2; C. J. Fewster, Lectures on quantum energy inequalities, arXiv:1208.5399., Secs. 2.3, 3.2.
\end{fact}

The same requirement can be written for the modes on a single light ray of a horizon. There it holds, on both orderings, exactly when the two-point kernel differs from the null vacuum's by a smooth kernel.

That gives the local form of (N$_{+}$) at every point, with no stationarity and no temperature assumed. It does not give the ground state itself: states thermal for the affine translations pass it everywhere, so a global datum such as boost stationarity at the edge is still needed.

\begin{theorem}[On a horizon generator the generator analogue of the microlocal (Hadamard) condition, on both orderings of the two-point kernel, is equivalent to a smooth difference from the null vacuum, so with no stationarity and no temperature it gives the local form of (N$_{+}$) at every point (with the stored theorem and boost stationarity at the edge, (N$_{+}$)); the affine-KMS states satisfy it, so it does not select the ground state]\label{c:synthesis.microlocal-condition-on-a-generator-gives-the-local-null-ground-state}
Physicists usually require that any reasonable quantum state, looked at closely enough, resembles empty space. The precise version of this requirement says that the very fine-grained structure of the state carries only positive energy. This result writes a version of that requirement for the quantum modes running along a single light ray of a horizon and shows that it is the same as saying the state differs from the vacuum only by something smooth. From there the store's earlier results give the horizon temperature at every point, without assuming that anything is static. It does not single out the vacuum, though: a gas of modes at any temperature also passes, so to get the full premise one still needs the state to be static near the point where the two halves of the horizon meet. Whether the series' states obey the requirement is still an open assumption.
\end{theorem}
\noindent\emph{Rests on:} \ref{c:synthesis.local-normality-at-the-edge-gives-the-null-ground-state}, \ref{c:synthesis.stationary-horizon-state-is-its-own-edge-scaling-limit}, \ref{c:synthesis.positive-null-energy-gives-r-premises}, \ref{c:synthesis.imported-microlocal-hadamard-condition}, \ref{c:synthesis.imported-local-quasi-equivalence-of-hadamard-states}, \ref{c:synthesis.imported-kay-wald-uniqueness-and-no-go}, \ref{c:synthesis.imported-propagation-of-hadamard-form}, \ref{c:synthesis.imported-powers-stormer-takesaki-disjointness}.

The next question is whether states of the field in the bulk obey this condition on the light rays. On the flat gravitational field the Dirac vacuum of every mass restricts on every null plane to the null vacuum, and every gauge-invariant quasi-free Hadamard Dirac state restricts to one satisfying the generator condition.

So, without gravity, the local form of (N$_{+}$) holds at every point of every local Rindler horizon; the ground-state clause, and fields of nonzero field strength, are not reached here.

\begin{theorem}[On the flat gravitational field every gauge-invariant quasi-free Dirac state differing from a vacuum by a smooth kernel restricts on every null plane to a state satisfying the generator microlocal condition, the vacuum itself to the null vacuum, so it has the local form of (N$_{+}$) on every local Rindler horizon; not the ground-state clause, not nonzero field strength]\label{c:synthesis.bulk-hadamard-states-restrict-to-the-generator-condition-on-flat-null-planes}
The series' account of horizon temperature rests on an assumption that, very close to any point of a horizon, the quantum field looks like empty space. This result checks that assumption against the field theory in empty flat space. On any light-like plane, the field's vacuum, whatever the particle's mass, looks exactly like the empty state that the assumption names, and every state that physicists regard as reasonable, one that differs from the vacuum only smoothly, looks like it at short distances too. So in flat space the assumption's short-distance part is a consequence of the field theory, though not the part that says the state is the vacuum everywhere, and not yet in a gravitational field.
\end{theorem}
\noindent\emph{Rests on:} \ref{c:synthesis.light-ray-good-part-of-a-spinor-is-two-chirality-arrays}, \ref{c:synthesis.slot-module-light-ray-arrays-are-weyl-components}, \ref{c:synthesis.positive-null-energy-gives-r-premises}, \ref{c:field-theory.second-quantisation}, \ref{c:synthesis.microlocal-condition-on-a-generator-gives-the-local-null-ground-state}, \ref{c:synthesis.imported-microlocal-hadamard-condition}, \ref{c:synthesis.imported-null-plane-restriction-of-free-fields-is-mass-independent}.

Where gravity is present, the first step is to look at a state very close to a single point. Every gauge-invariant quasi-free Hadamard state of the Dirac field of any mass, at every point of a globally hyperbolic gravitational field, scales to the flat massless vacuum on both orderings.

Under the scaling the mass term fades as $O(l)$, the connection and transport terms together as $O(l^{2})$, and the remainder as $O(l^{3})$. This is a standard theorem of Sahlmann and Verch, re-derived in the spacetime algebra for Dirac spinors, given the imported Hadamard parametrix. By itself it says nothing about the state restricted to the light rays of a horizon.

\begin{theorem}[Given the Hadamard parametrix of the squared Dirac operator (imported), every gauge-invariant quasi-free Hadamard state of the Dirac field of mass m0 >= 0 on a globally hyperbolic gravitational field scales, under the store's dilation in a gauge with hb(x0) = 1, to the flat field's massless vacuum on both orderings (standard: Sahlmann and Verch 2001, Prop. 5.6, re-derived in the spacetime algebra for Dirac spinors), so clause (a) of (HS) holds and the premise's content is its generator clause (b)]\label{c:synthesis.hadamard-dirac-states-scale-to-the-flat-massless-vacuum-at-every-point-of-a-gravitational-field}
The series assumed that any state of the free electron field of the standard regular kind (gauge-invariant, quasi-free and Hadamard), looked at closely enough at a point where gravity is present, resembles the empty state of flat space. This result proves that assumption for the bulk of space-time, for massive and massless particles and for every gravitational field with a well-posed causal structure (globally hyperbolic), by redoing in the series' algebra a known theorem of Sahlmann and Verch from the mathematical physics literature: close up, the curvature, the mass and the gravitational connection all fade out of the state's correlations at definite rates. It relies on two standard analytic facts quoted from the literature rather than proved in the series. What it does not prove is the companion statement the series also needs: that the same is true when the state is looked at only along the light rays of a small horizon. That part stays an assumption.
\end{theorem}
\noindent\emph{Rests on:} \ref{c:synthesis.imported-gtg-dictionary}, \ref{c:synthesis.imported-dimock-dirac-car-algebra-on-a-gravitational-field}, \ref{c:field-theory.second-quantisation}, \ref{c:synthesis.imported-microlocal-hadamard-condition}, \ref{c:synthesis.imported-dirac-hadamard-form-propagates-and-adiabatic-states-are-hadamard}, \ref{c:synthesis.imported-hadamard-parametrix-of-the-squared-dirac-operator-leads-with-van-vleck-root-times-spin-transport}, \ref{c:synthesis.imported-free-field-spin-statistics-and-guido-longo-theorem}, \ref{c:synthesis.bulk-hadamard-states-restrict-to-the-generator-condition-on-flat-null-planes}.

Looking only along the light rays might let the tidal part of gravity into the state at short distances. On any null hypersurface generated by null geodesics, the Weyl and Ricci field strength along the generators, the expansion and the shear all cancel from the coefficient of the leading null singularity, because $\Delta^{1/2}/\sqrt{\det M} = (a(u)a(u'))^{-1/2}$.

Since Hadamard states differ there by smooth kernels, the question becomes one about the geometry of the field: it holds for every Hadamard state or for none.

\begin{theorem}[On every null hypersurface of a gravitational field generated by affine null geodesics, the field strength along the generators, the expansion and the shear cancel exactly from the coefficient of the leading null singularity of a Hadamard Dirac kernel restricted to it, which at leading order in the canonical normalisation is the null vacuum; restricted kernels of Hadamard states differ by smooth kernels, so clause (b) of (HS) is a property of the field and reduces to a remainder bound on the Hadamard parametrix, argued but not proved]\label{c:synthesis.field-strength-cancels-from-the-leading-null-singularity-of-hadamard-dirac-states-on-null-hypersurfaces}
The series assumes that, looked at closely along the light rays that make up a small horizon, any reasonable quantum state of the electron field resembles empty space, even where gravity is present. One way this could fail is if gravity's tidal stretching of the light rays left a mark on the state at the shortest distances. This result shows that, at the leading order, it does not: the tidal effects change how light rays spread and how distances between them are measured, but the two changes cancel exactly. It also shows that the assumption is either true for all reasonable states or for none, so it is really a question about gravity's geometry, not about which state the world is in. Smaller correction terms needed a technical estimate, argued here and proved by a later result.
\end{theorem}
\noindent\emph{Rests on:} \ref{c:synthesis.imported-gtg-dictionary}, \ref{c:synthesis.imported-microlocal-hadamard-condition}, \ref{c:synthesis.imported-dirac-hadamard-form-propagates-and-adiabatic-states-are-hadamard}, \ref{c:synthesis.imported-hadamard-parametrix-of-the-squared-dirac-operator-leads-with-van-vleck-root-times-spin-transport}, \ref{c:synthesis.imported-dimock-dirac-car-algebra-on-a-gravitational-field}, \ref{c:synthesis.bulk-hadamard-states-restrict-to-the-generator-condition-on-flat-null-planes}, \ref{c:synthesis.microlocal-condition-on-a-generator-gives-the-local-null-ground-state}.

The field-strength result left one estimate on the smaller terms argued rather than proved, and the next result proves it. On every piece of a geodesic null hypersurface whose distinct generators are spacelike separated, every local Rindler horizon included, the Dirac Hadamard parametrix restricted to the generators is $W_{0}\otimes\mathrm{id}_{K}$ plus a bounded kernel, for every mass.

The two logarithmically growing pieces cancel identically, the expansion term against the projected transverse spin connection, so the bound that (HS) posed holds. This is shown for one way of restricting to the light rays, a regularised limit; that it agrees with the restriction that applies to the states themselves is not part of it.

\begin{theorem}[On every piece of a geodesic null hypersurface of a gravitational field whose distinct generators are spacelike separated (every small enough piece, every local Rindler horizon), the Dirac Hadamard parametrix, restricted to the generators in the canonical normalisation as the epsilon-regularised transverse-smeared limit, is the null vacuum plus a bounded kernel: the log epsilon terms cancel identically, so the remainder bound posed by (HS) holds for every mass]\label{c:synthesis.restricted-dirac-hadamard-parametrix-on-a-null-hypersurface-is-the-null-vacuum-plus-a-bounded-kernel}
The series needs every reasonable quantum state of the electron field to look like empty space along the light rays of the small horizons used to derive the equations of gravity, when looked at closely enough. The question had been reduced to a statement about gravity's geometry alone: a fixed mathematical object, which carries the short-distance behaviour shared by all reasonable states, should look like that of empty space along those light rays, up to corrections that fade at short distances. The strongest part had been checked before; two weaker pieces, which separately would grow without limit, were expected but not shown to cancel. This result shows that they do cancel exactly, on small enough pieces of the light sheets of every gravitational field (including the small horizons the series uses) and for every mass of the field, and that everything left over stays bounded. So the assumed statement is now proved, for the precise way the object is restricted to the light rays. One step remains before it applies to the states themselves: showing that this way of restricting agrees with the more fundamental one.
\end{theorem}
\noindent\emph{Rests on:} \ref{c:synthesis.imported-gtg-dictionary}, \ref{c:synthesis.imported-hadamard-parametrix-of-the-squared-dirac-operator-leads-with-van-vleck-root-times-spin-transport}, \ref{c:synthesis.imported-dirac-hadamard-form-propagates-and-adiabatic-states-are-hadamard}, \ref{c:synthesis.field-strength-cancels-from-the-leading-null-singularity-of-hadamard-dirac-states-on-null-hypersurfaces}, \ref{c:synthesis.bulk-hadamard-states-restrict-to-the-generator-condition-on-flat-null-planes}.

Restricting a state itself to the light rays needs a precise meaning, and the field's own current supplies one. Because the Dirac current is future causal, the trace on a null hypersurface is a contraction, so every Dirac state restricts to the generators, with $W^{+} + W^{-} = \iota^{*}\iota \le 1$.

Restricted kernels of Hadamard states differ by smooth pullbacks, with no further condition. If the fermion anticommutation relations hold on the generators, the restriction is also the limit of the bulk two-point function over thinner and thinner layers.

\begin{theorem}[Every state of the Dirac field on a gravitational field restricts to a null hypersurface through the flux contraction of the conserved current, with no characteristic Cauchy problem, its restricted kernels have generator scaling degree at most one, and restricted kernels of Hadamard states differ by smooth pullbacks; if the generator anticommutation relations hold on nearby null hypersurfaces (proved in the literature for light cones), the restriction is the thin-layer limit of the bulk two-point function, and the restricted Hadamard parametrix then has a regularisation-free definition, whose agreement with the regularised one used in (HS) is argued]\label{c:synthesis.dirac-states-restrict-to-null-hypersurfaces-through-the-flux-contraction-and-are-thin-layer-limits-under-the-generator-car}
The series assumes that, seen closely along the light rays of a small horizon, quantum states of the electron field look like empty space. Testing this needs a precise way to look at a state "on the light rays only", and the earlier work borrowed one from memory. This result builds it from a basic fact: the flow of electron charge-current is never faster than light, so what flows through a sheet of light rays can never exceed the total. That alone makes every state restrict to the light rays, caps how singular the restriction can be, and shows that any two reasonable states differ there only by a smooth amount. If the light-ray field also obeys the usual rules for fermions, which mathematicians have proved for light cones but which are only believed for general light sheets, then the restriction is simply the limit of ordinary measurements made in thinner and thinner slabs, and the earlier borrowed step is no longer needed. Whether this matches the way the horizon assumption itself is computed is argued, not proved.
\end{theorem}
\noindent\emph{Rests on:} \ref{c:synthesis.imported-gtg-dictionary}, \ref{c:synthesis.imported-dimock-dirac-car-algebra-on-a-gravitational-field}, \ref{c:synthesis.imported-microlocal-hadamard-condition}, \ref{c:synthesis.imported-dirac-hadamard-form-propagates-and-adiabatic-states-are-hadamard}, \ref{c:synthesis.imported-hadamard-parametrix-of-the-squared-dirac-operator-leads-with-van-vleck-root-times-spin-transport}, \ref{c:synthesis.imported-haefner-nicolas-dirac-characteristic-problem-on-light-cones-near-the-vertex}, \ref{c:synthesis.field-strength-cancels-from-the-leading-null-singularity-of-hadamard-dirac-states-on-null-hypersurfaces}, \ref{c:synthesis.bulk-hadamard-states-restrict-to-the-generator-condition-on-flat-null-planes}, \ref{c:synthesis.microlocal-condition-on-a-generator-gives-the-local-null-ground-state}.

That remaining condition is supplied next. On every geodesic null hypersurface the Dirac causal propagator, taken between the hypersurface and a nearby null leaf, tends to $\delta(u - u')\delta^{2}(y - y')\Lambda$, so the generator anticommutation relations hold, local Rindler horizons included.

The correction that could have spoiled this vanishes by the same half-density law that removed the field strength.

\begin{theorem}[On every geodesic null hypersurface of a gravitational field the causal propagator of the Dirac field, restricted between the hypersurface and a nearby leaf of a null foliation, tends to the canonical delta kernel of the generators, the step term cancelling by the half-density law; with the flux contraction of the conserved current, the generator anticommutation relations hold there, and the restricted Hadamard parametrix in the thin-layer definition is the null vacuum plus a bounded kernel]\label{c:synthesis.dirac-causal-propagator-between-nearby-null-leaves-tends-to-the-generator-car-on-every-geodesic-null-hypersurface}
Earlier work needed a basic rule for the electron field seen only along the light rays of a small horizon: the field there should obey the standard rules for fermions, exactly as in empty space. Mathematicians had proved this for light cones; for the general light sheets used in the series it was only believed. This result proves it for every such sheet by a short calculation with the one object every state shares, the propagator that carries signals at or below light speed. Two corrections that could have spoiled the rule turn out to cancel, because the area of the sheet and the focusing of light rays are tied together. Together with two companion results, it also shows that every reasonable state looks like empty space along these light rays at short distances.
\end{theorem}
\noindent\emph{Rests on:} \ref{c:synthesis.imported-gtg-dictionary}, \ref{c:synthesis.imported-dimock-dirac-car-algebra-on-a-gravitational-field}, \ref{c:synthesis.imported-hadamard-parametrix-of-the-squared-dirac-operator-leads-with-van-vleck-root-times-spin-transport}, \ref{c:synthesis.imported-dirac-hadamard-form-propagates-and-adiabatic-states-are-hadamard}, \ref{c:synthesis.imported-microlocal-hadamard-condition}, \ref{c:synthesis.field-strength-cancels-from-the-leading-null-singularity-of-hadamard-dirac-states-on-null-hypersurfaces}, \ref{c:synthesis.bulk-hadamard-states-restrict-to-the-generator-condition-on-flat-null-planes}, \ref{c:synthesis.dirac-states-restrict-to-null-hypersurfaces-through-the-flux-contraction-and-are-thin-layer-limits-under-the-generator-car}, \ref{c:synthesis.restricted-dirac-hadamard-parametrix-on-a-null-hypersurface-is-the-null-vacuum-plus-a-bounded-kernel}.

The pieces now assemble into one theorem. Theorem (HS), formerly a premise: for every gauge-invariant quasi-free Hadamard state of the Dirac field of any mass, at every point of every local Rindler horizon, the restriction to the generators is $W_{0}\otimes\mathrm{id}_{K}$ plus a bounded and a smooth kernel, so its generator scaling limit is the null vacuum.

Hence the local form of (N$_{+}$) holds there for every such state. A few standard inputs, among them the local null foliation and the Morse lemma with parameters, are cited from memory rather than imported with a source. The theorem gives neither the ground-state clause of (N$_{+}$) nor which state the world is in.

\begin{theorem}[(HS), the Hadamard scaling limit, a theorem (formerly a premise): for every gauge-invariant quasi-free Hadamard state of the Dirac field of any mass on a gravitational field, at every point of every local Rindler horizon, the bulk scaling limit is the flat massless vacuum and the restriction to the generators through the flux trace has the null vacuum as its generator scaling limit, the epsilon-regularised restriction of the parametrix agreeing with it modulo a bounded time-function-dependent kernel; hence the local form of (N$_{+}$) at every such point for these states]\label{c:synthesis.hadamard-scaling-limit-premise}
Physicists require that any reasonable quantum state, looked at closely enough, resembles empty space. The series had assumed this along the light rays of the small horizons used to derive the equations of gravity, wherever gravity is present. It is now proved, for the electron field of any mass, given standard results from the literature: the flow of electric charge through a sheet of light can never exceed the total, signals travel at or below light speed, and two growing corrections cancel because the area of the sheet and the focusing of its light rays are tied together. So every reasonable state looks like empty space along those light rays at short distances, which is what the series' horizon temperature needs locally. It does not pick out which state the world is in, and it does not by itself give the full horizon assumption, since warm states pass the test too.
\end{theorem}
\noindent\emph{Rests on:} \ref{c:synthesis.imported-dimock-dirac-car-algebra-on-a-gravitational-field}, \ref{c:synthesis.imported-microlocal-hadamard-condition}, \ref{c:synthesis.imported-dirac-hadamard-form-propagates-and-adiabatic-states-are-hadamard}, \ref{c:synthesis.imported-hadamard-parametrix-of-the-squared-dirac-operator-leads-with-van-vleck-root-times-spin-transport}, \ref{c:synthesis.hadamard-dirac-states-scale-to-the-flat-massless-vacuum-at-every-point-of-a-gravitational-field}, \ref{c:synthesis.field-strength-cancels-from-the-leading-null-singularity-of-hadamard-dirac-states-on-null-hypersurfaces}, \ref{c:synthesis.restricted-dirac-hadamard-parametrix-on-a-null-hypersurface-is-the-null-vacuum-plus-a-bounded-kernel}, \ref{c:synthesis.dirac-states-restrict-to-null-hypersurfaces-through-the-flux-contraction-and-are-thin-layer-limits-under-the-generator-car}, \ref{c:synthesis.dirac-causal-propagator-between-nearby-null-leaves-tends-to-the-generator-car-on-every-geodesic-null-hypersurface}, \ref{c:synthesis.bulk-hadamard-states-restrict-to-the-generator-condition-on-flat-null-planes}.

That leaves the global clause of the null ground-state premise, the part no Hadamard condition reaches. For the fermionic modes of a horizon generator, a state whose restriction beyond each point of a set unbounded below is KMS, at any period, for the boost about that point is the null vacuum of (N$_{+}$), with every period forced to $2\pi$.

Two points already fix the period, because dilations about two points generate translations and the thermal kernel $\sinh(s/2)/\sinh(\pi s/\beta)$ is translation invariant only at $\beta = 2\pi$. One point, or a set of points bounded below, does not suffice.

This is an equivalent form of the ground-state clause, not a new premise; it applies on Minkowski null planes and, argued, on de Sitter horizons, but not on the eternal hole.

\begin{theorem}[For the fermionic (CAR) modes of a horizon generator, a state whose restriction to the future half beyond each point of a set unbounded below is KMS, at some period, for the boost about that point is the null vacuum of (N$_{+}$), and every period is $2\pi$: two points already fix the period, so (N$_{+}$) on a generator is equivalent to temperature-free thermality of its future halves, with no stationarity or quasi-freeness assumed; it applies on generators each of whose points is the edge of a Killing wedge (Minkowski null planes; de Sitter horizons, argued), not on the eternal hole]\label{c:synthesis.thermality-of-every-future-half-generator-fixes-the-period-and-the-null-ground-state}
The series assumes that the quantum modes running along a horizon's light rays are in their lowest-energy state. For particles of the electron's kind (fermions), this result gives an equivalent condition that sounds quite different: from every point along a light ray, looking forward along the ray, the modes appear warm, as the modes beyond an accelerated observer's horizon do. Surprisingly, the temperature does not need to be specified: asking for warmth about two different points of the same ray already fixes it at the standard value, and asking for it about points reaching arbitrarily far back fixes the whole state as the lowest-energy one. This links the assumption to a well-studied property of the vacuum in relativistic quantum theory, and it shows how it could hold on the horizon of an expanding universe for massive particles too, given an assumption about how the horizon's modes relate to the region behind it, though not, by this route, on a black hole that has always existed.
\end{theorem}
\noindent\emph{Rests on:} \ref{c:hawking.i-state-continuum}, \ref{c:hawking.imported-ideal-fermi-gas-uniqueness}, \ref{c:synthesis.positive-null-energy-gives-r-premises}, \ref{c:synthesis.stationary-horizon-state-is-its-own-edge-scaling-limit}, \ref{c:synthesis.de-sitter-horizon-translation-is-conformal-null-rotation}, \ref{c:angle.imported-theorems-deliver-their-hypotheses}.

\section{Where the state comes from}

This section turns from the horizon to the state itself: where the series' states come from, and whether they are of the well-behaved kind the horizon results need.

Everything so far assumes that the series' states are Hadamard in the first place. The premise (HD) adopts this: the state the descent deposits is a Hadamard state on a Cauchy collar of the descent, for the Dirac and scalar fields, with its occupation profile left free.

Because the Hadamard form propagates, every later field then carries Hadamard states, and with (HS) the local form of (N$_{+}$) follows at every horizon. The premise gives neither which state, nor the ground-state clause, nor the spectrum of primordial ripples.

The first test is the kind of deposit the series uses, a thermal one. On a descended field of Robertson-Walker form with smooth scale factor (a stated hypothesis, as for a descent of radiation type), a homogeneous massless Dirac deposit diagonal in the conformal vacuum's modes is Hadamard exactly when every moment of its occupation profile is finite.

So the thermal deposit is Hadamard. The number of quanta and the spread of their energies do not decide it: a $(1 + p/T)^{-6}$ tail has both finite and is not Hadamard.

\begin{theorem}[On a descended field of Robertson-Walker form with smooth scale factor (as for a descent of radiation type, a stated hypothesis), a homogeneous deposit of the massless Dirac field diagonal in the conformal vacuum's quasi-particle modes (no pair coherence) is Hadamard exactly when every moment of its occupation profile is finite, the standard flat criterion carried by conformal covariance: the thermal deposit is Hadamard, while the number and energy spread of the quanta, the data of the stored shot-noise bound, do not decide it]\label{c:synthesis.deposit-over-the-conformal-vacuum-is-hadamard-exactly-when-every-occupation-moment-is-finite}
When the universe's earlier phase handed its contents to the later one, it left the quantum fields in some state. Physics in a gravitational field needs that state to be of a particular well-behaved kind, called Hadamard, which roughly means it looks like empty space at very short distances in a smooth way. This result shows, for massless spin-one-half particles and assuming the young universe was uniform and expanding smoothly, that a thermal hand-over, the kind the series uses, is of that kind. It also shows that the series' earlier calculation of how much the hand-over can disturb things, which uses only how many particles there are and how their energies spread, cannot settle the question: what matters is how fast the number of particles falls off at high energy.
\end{theorem}
\noindent\emph{Rests on:} \ref{c:synthesis.quasi-free-deposit-is-vacuum-plus-shot-noise}, \ref{c:synthesis.imported-microlocal-hadamard-condition}, \ref{c:hawking.i-state}, \ref{c:field-theory.second-quantisation}, \ref{c:synthesis.imported-gtg-dictionary}, \ref{c:synthesis.weyl-flat-fields-need-no-scaling-limit-premise-and-robertson-walker-generators-see-the-schwarzian}, \ref{c:synthesis.imported-conformal-covariance-of-massless-fields-and-local-conformal-flatness}, \ref{c:synthesis.imported-bochner-theorem-moment-lemma-and-whitney-even-functions}.

Massive fields are less forgiving. For a massive Dirac field on a Robertson-Walker field, at an instant where $ma' \neq 0$ the instantaneous vacuum differs from an adiabatic state of infinite order by $|\beta_{k}|^{2} = M'^{2}/16k^{4}$, a finite number of quanta with logarithmically divergent energy, so it is not Hadamard, and nor is any deposit diagonal over it.

A descent across which a derivative of $M = ma$ jumps leaves no Hadamard state either, so a Hadamard deposit needs a smooth descent.

\begin{theorem}[On a spatially flat Robertson-Walker field, for a massive Dirac field, at an instant where $ma' \neq 0$ the instantaneous vacuum differs from an adiabatic state of infinite order by $|\beta_{k}|^{2} = M'^{2}/16k^{4}$ (finite number, logarithmically divergent energy) and is not Hadamard, nor is any deposit diagonal over it; a state adiabatic of infinite order before a descent across which the $n$-th derivative of $M = ma$ jumps has $|\beta_{k}| = |\Delta M^{(n)}|/(2k)^{n+1}$ after it and is not Hadamard for any finite $n$]\label{c:synthesis.massive-dirac-instantaneous-vacuum-and-non-smooth-descents-are-not-hadamard}
The series needs the quantum state left behind by the descent to be of a regular kind, called Hadamard, so that energies can be computed. An earlier result described that state as a vacuum plus a sprinkling of particles, without saying which vacuum. For electrons and their kin, on an expanding universe, the vacuum is not unique: different recipes give different answers. This result tests the most naive recipe, taking the lowest-energy state at a single instant, and finds that it fails: compared with a regular state it carries a thin spray of extra particles at high energies, few in number but with infinite energy, and adding particles on top of it does not repair this. The same calculation shows that any abrupt descent fails too: a sudden change, even only in how fast the universe grows, leaves a high-energy tail that spoils the regularity. A simpler model stored earlier, with a sudden start, fails in the same way.
\end{theorem}
\noindent\emph{Rests on:} \ref{c:synthesis.imported-microlocal-hadamard-condition}, \ref{c:synthesis.sudden-descent-feature-not-tilt}, \ref{c:synthesis.imported-dirac-hadamard-form-propagates-and-adiabatic-states-are-hadamard}, \ref{c:synthesis.imported-ford-parker-infrared-divergence}.

If smoothness does not pick the state, something else must, and a natural way to select it from the Dirac field's own structure is its spin product. Split against the conserved current over a collar, the spin product is odd under the chirality $\Gamma\psi = I_{4}\psi I_{4}$ for massless fields, so the state it selects has occupation one half over the conformal vacuum on every spatially flat Robertson-Walker collar and is Hadamard nowhere.

For massive fields it is Hadamard only with a smooth weight, a choice the construction does not make. So the spin product does not select the deposit that (HD) assumes, while the twistor form selects the conformal vacuum.

\begin{theorem}[The state selected by the Dirac field's own spin product (the Fewster-Lang weighted spin-product state on a collar, the finite-time form of the fermionic projector state) is odd under chirality for the massless field and so is Hadamard on no open set of any collar on which its operator is injective, with occupation one half over the conformal vacuum on Robertson-Walker collars for every weight; for massive fields it is Hadamard only with a smooth time weight; so the spin product does not select the deposit that (HD) assumes]\label{c:synthesis.spin-product-state-of-the-massless-dirac-field-is-odd-under-chirality-and-never-hadamard}
Quantum fields need a well-behaved state, one that looks like empty space at very short distances. The series' account of the earliest universe leaves its fields in some state but does not say which, so it is assumed to be well behaved. One natural way to pick a state using only the electron field's own structure is to split its solutions with a mixed-sign quantity built from the field, the one that gives an electron its mass term, averaged over a slice of time. This result shows that for massless particles, which is how all particles behave at the universe's start, that split always picks a badly behaved state: the quantity changes sign under the swap of left- and right-handed particles, which forces the chosen state to be half filled at every energy. For massive particles it works only if the time slice is given soft edges, a choice the method does not make. A different mixed-sign quantity, from the series' twistor picture, does pick a well-behaved state.
\end{theorem}
\noindent\emph{Rests on:} \ref{c:synthesis.imported-dimock-dirac-car-algebra-on-a-gravitational-field}, \ref{c:synthesis.imported-microlocal-hadamard-condition}, \ref{c:synthesis.massive-dirac-instantaneous-vacuum-and-non-smooth-descents-are-not-hadamard}, \ref{c:synthesis.deposit-over-the-conformal-vacuum-is-hadamard-exactly-when-every-occupation-moment-is-finite}, \ref{c:synthesis.imported-conformal-covariance-of-massless-fields-and-local-conformal-flatness}, \ref{c:synthesis.attachment-conformal-energy-is-generated-by-z}, \ref{c:synthesis.imported-fewster-lang-spin-product-states-of-the-dirac-field}.

\section{What the Hadamard condition leaves free}

The last section asks what the Hadamard condition leaves undecided about the deposit.

One might hope that the Hadamard condition also fixes the pattern of primordial ripples. It does not: relative to any homogeneous Hadamard reference, a homogeneous deposit can have any spectrum on any bounded band of scales.

What it forbids is a red law on a band reaching $k = 0$ for slice fields of nonzero mean, and, for the conformally coupled field, a power law on a band unbounded toward short scales. So a red band needs two ends, and the condition does not say where they are.

\begin{theorem}[Measured against a homogeneous Hadamard reference state, the Hadamard condition leaves a homogeneous deposit's spectrum free on every bounded band of scales; for the conformally coupled field it allows a power-law ratio on no band unbounded toward short scales, and, for slice fields of nonzero mean, a red law on no band reaching k = 0, so a red band needs an ultraviolet and an infrared end, neither of which the condition fixes]\label{c:synthesis.hadamard-condition-leaves-the-deposit-spectrum-free-on-bounded-bands}
The series needs the quantum state left by the descent to be of a regular kind, called Hadamard, so that energies and horizon physics make sense. One might hope that this regularity would also fix the pattern of primordial ripples seen in the sky. This result shows it does not. The Hadamard condition only controls the shortest distances; over any finite range of scales, including the whole range astronomers measure, the state can have any strength and any tilt of ripples. What the condition and the existence of the state do require is that a pattern tilted like the observed one must stop somewhere at very short scales and somewhere at very long ones, and nothing in the condition says where.
\end{theorem}
\noindent\emph{Rests on:} \ref{c:synthesis.imported-microlocal-hadamard-condition}, \ref{c:synthesis.imported-de-sitter-invariant-hadamard-state-unique}, \ref{c:synthesis.frozen-point-geometry-reads-descent-as-vacuum}, \ref{c:standard-model.imported-frozen-point-dictionary}, \ref{c:synthesis.imported-ford-parker-infrared-divergence}, \ref{c:synthesis.imported-conformal-covariance-of-massless-fields-and-local-conformal-flatness}.

\section{Results used from elsewhere}
The proofs above use the following results of the series, stated briefly; each is proved in the paper named.

\begin{fact}[Imported: the theorems that deliver $2\pi$, and their hypotheses]\label{c:angle.imported-theorems-deliver-their-hypotheses}
Under the spectrum condition, covariance and locality, the modular flow of a wedge or half-sided algebra is the boost at rate $2\pi$, and a regular stationary quasi-free state at a bifurcate Killing horizon is KMS at $\kappa/2\pi$.
\emph{Source:} J. J. Bisognano and E. H. Wichmann, On the duality condition for a Hermitian scalar field, J. Math. Phys. 16 (1975), 985--1007.; J. J. Bisognano and E. H. Wichmann, On the duality condition for quantum fields, J. Math. Phys. 17 (1976), 303--321.; H.-J. Borchers, The CPT-theorem in two-dimensional theories of local observables, Commun. Math. Phys. 143 (1992), 315--332.; P. D. Hislop and R. Longo, Modular structure of the local algebras associated with the free massless scalar field theory, Commun. Math. Phys. 84 (1982), 71--85.; R. Brunetti, D. Guido and R. Longo, Modular structure and duality in conformal quantum field theory, Commun. Math. Phys. 156 (1993), 201--219.; G. L. Sewell, Quantum fields on manifolds: PCT and gravitationally induced thermal states, Ann. Phys. 141 (1982), 201--224.; B. S. Kay and R. M. Wald, Theorems on the uniqueness and thermal properties of stationary, nonsingular, quasifree states on spacetimes with a bifurcate Killing horizon, Phys. Rep. 207 (1991), 49--136.; W. G. Unruh, Notes on black-hole evaporation, Phys. Rev. D 14 (1976), 870--892..
\end{fact}

\begin{fact}[Second quantisation]\label{c:field-theory.second-quantisation}
With the single-algebra field and $i = I_4$, second quantisation gives a local field and the standard $d\Gamma(M)$, one chirality entering complex-linearly and the other conjugated.
\emph{Source:} FT, Prop 5.2.
\end{fact}

\begin{fact}[Gibbs states are KMS, and conversely, in finite dimensions]\label{c:hawking.gibbs-states-are-kms-conversely-finite}
On a full matrix algebra with dynamics generated by a diagonal Hamiltonian, the Gibbs state at inverse temperature $\beta$ is KMS, and it is the unique KMS state at that $\beta$, degenerate spectrum included.
\emph{Source:} T7, Lem 3.1 (\href{https://doi.org/10.5281/zenodo.22722121}{doi:10.5281/zenodo.22722121}).
\end{fact}

\begin{fact}[Imported: the exterior is exactly general relativity]\label{c:hawking.imported-exterior-exactly-general-relativity}
Where the gauge field is invertible and $\Lambda = 0$, the series' description coincides exactly with general relativity, its vacuum equations being the parallelism of the scale tractor.
\emph{Source:} T7, Fact 2.5 (\href{https://doi.org/10.5281/zenodo.22722121}{doi:10.5281/zenodo.22722121}).
\end{fact}

\begin{fact}[Imported: ideal-Fermi-gas uniqueness]\label{c:hawking.imported-ideal-fermi-gas-uniqueness}
For free fermions with quasi-free dynamics the $\beta$-KMS state is unique and is the quasi-free state with Fermi-Dirac symbol $(1 + e^{\beta h})^{-1}$.
\emph{Source:} O. Bratteli and D. W. Robinson, Operator Algebras and Quantum Statistical Mechanics 2, 2nd ed., Springer, Berlin (1997)., Ex.~5.3.2 and \S 5.3; H. Araki, On quasifree states of CAR and Bogoliubov automorphisms, Publ. RIMS Kyoto Univ. 6 (1971), 385--442..
\end{fact}

\begin{fact}[Imported: KMS passes to inductive limits]\label{c:hawking.imported-kms-passes-inductive-limits}
A state that is KMS at $\beta$ on each subalgebra of an increasing chain of dynamics-invariant subalgebras is KMS at $\beta$ on the closure of their union.
\emph{Source:} O. Bratteli and D. W. Robinson, Operator Algebras and Quantum Statistical Mechanics 2, 2nd ed., Springer, Berlin (1997)., Prop.~5.3.25.
\end{fact}

\begin{fact}[Imported: modular descent, its premises, and the geometric normalisation]\label{c:hawking.imported-modular-descent-premises-geometric-normalisation}
A pure, boost-invariant, entangling state of a horizon pair descends to a KMS state at an undetermined temperature, which the premises (P) and (R) fix as Fermi-Dirac at the Hawking temperature $\TH = 1/8\pi M$.
\emph{Source:} T7, Fact 2.2 (\href{https://doi.org/10.5281/zenodo.22722121}{doi:10.5281/zenodo.22722121}).
\end{fact}

\begin{fact}[Imported: quasi-free KMS reduces to the two-point level]\label{c:hawking.imported-quasi-free-kms-reduces-two}
A gauge-invariant quasi-free CAR state is KMS for quasi-free dynamics exactly when its two-point function satisfies the two-point KMS identity.
\emph{Source:} H. Araki, On quasifree states of CAR and Bogoliubov automorphisms, Publ. RIMS Kyoto Univ. 6 (1971), 385--442.; O. Bratteli and D. W. Robinson, Operator Algebras and Quantum Statistical Mechanics 2, 2nd ed., Springer, Berlin (1997)., \S 5.3.
\end{fact}

\begin{fact}[Imported: quasi-free states of the CAR algebra]\label{c:hawking.imported-quasi-free-states-car-algebra}
Each symbol $0 \le A \le 1$ determines a unique gauge-invariant quasi-free state of the CAR algebra; quasi-free states are fixed by their symbol and pairing, and gauge invariance is the vanishing of the pairing.
\emph{Source:} H. Araki, On quasifree states of CAR and Bogoliubov automorphisms, Publ. RIMS Kyoto Univ. 6 (1971), 385--442.; O. Bratteli and D. W. Robinson, Operator Algebras and Quantum Statistical Mechanics 2, 2nd ed., Springer, Berlin (1997)., \S 5.2.
\end{fact}

\begin{fact}[Imported: the s-wave bosonic sector]\label{c:hawking.imported-s-wave-bosonic-sector}
Boundary-mode currents satisfy the free-boson relations exactly up to an edge term, and under Premise (Q) they carry Bose-Einstein occupancy at $\TH$ in the s-wave sector.
\emph{Source:} T7, Fact 2.4 (\href{https://doi.org/10.5281/zenodo.22722121}{doi:10.5281/zenodo.22722121}).
\end{fact}

\begin{fact}[Ascended marginals are diagonal]\label{c:singularity.ascended-marginals-are-diagonal}
Under univalence every single-mode marginal is diagonal in the number basis, hence equal to its descent and, at nonzero weight, a Gibbs state of the boost generator at a unique $\beta_k$.
\emph{Source:} T5, Thm 7.10 (\href{https://doi.org/10.5281/zenodo.22721790}{doi:10.5281/zenodo.22721790}).
\end{fact}

\begin{fact}[The curvature of the gauge field: exact shape decoupling]\label{c:singularity.curvature-gauge-field-exact-shape-decoupling}
The curvature of the black-hole gauge field is exactly $\mathcal{R} = -(M/r^3)\,S$ with $S = 1 - 3\Pi$ a fixed shape operator of spectrum $(-2,-2,1,1,1,1)$.
\emph{Source:} T5, Thm 6.1 (\href{https://doi.org/10.5281/zenodo.22721790}{doi:10.5281/zenodo.22721790}).
\end{fact}

\begin{fact}[The descended form is state independent]\label{c:singularity.descended-form-state-independent}
Whatever the ascended state, each descended single-mode marginal is a Gibbs state of the boost generator at a unique $\beta_k$, for every mode of nonzero weight.
\emph{Source:} T5, Thm 7.6 (\href{https://doi.org/10.5281/zenodo.22721790}{doi:10.5281/zenodo.22721790}).
\end{fact}

\begin{fact}[Imported: surface gravity as redshifted acceleration]\label{c:singularity.imported-surface-gravity-redshifted-acceleration}
For a static Killing horizon, the surface gravity, $(K \cdot D)K = \kappa K$ on the horizon locus, equals the horizon limit of the redshift-weighted proper acceleration of static observers.
\emph{Source:} R. M. Wald, General Relativity, University of Chicago Press (1984), §12.5.
\end{fact}

\begin{fact}[KMS uniqueness at a factor]\label{c:singularity.kms-uniqueness-factor}
On a matrix factor, for each nonzero inverse temperature the only KMS state of an inner flow is the Gibbs state of its generator.
\emph{Source:} T5, Thm 7.13 (\href{https://doi.org/10.5281/zenodo.22721790}{doi:10.5281/zenodo.22721790}).
\end{fact}

\begin{fact}[Modular structure of the completed pair]\label{c:singularity.modular-structure-completed-pair}
The element $\sigma = v^*v - vv^*$ of the completed pair is at once the boost generator, the shifted number operator and the generator of the descent projectors, and its boost flow is, up to reparametrisation, the only real rotor flow of the atom.
\emph{Source:} T5, Thm 6.3 (\href{https://doi.org/10.5281/zenodo.22721790}{doi:10.5281/zenodo.22721790}).
\end{fact}

\begin{fact}[Imported: the frozen point and the dictionary]\label{c:standard-model.imported-frozen-point-dictionary}
Mass is conformal weight read against the frozen point, $m^{2} = \Delta(3-\Delta)H^{2}$, so heavy fields have complex weight and the Planck mass's imaginary weight is $N \approx M_{P}/H$, about $10^{61}$.
\emph{Source:} T12, Fact 2.5 (\href{https://doi.org/10.5281/zenodo.22722669}{doi:10.5281/zenodo.22722669}).
\end{fact}

\begin{fact}[The twistor form on the conformal spinor space is the grading of the central rotation $Z = e_{0}\nu$ against the Hilbert form of the ladder-pair module, so the twistor attachment's conformal energy $|s| + 1 + n$ is the nonnegative generator of $Z$; every future translation, null special conformal transformation and de Sitter horizon generator is nonnegative on its modules, and the attachment's Fock vacuum satisfies the positivity and invariance clause of (N$_{+}$) on the local horizons of Minkowski space (and of de Sitter space for conformally invariant fields); on those horizons (N$_{+}$) is the choice of that vacuum as the state of the light rays, a restatement of the premise and not a derivation of it]\label{c:synthesis.attachment-conformal-energy-is-generated-by-z}
On the conformal spinor space the invariant Hermitian form is, up to a real factor, $\ip{\psi}{iZ\phi}$ with $Z = e_{0}\nu$, so $Z$ grades the twistor form into the Hilbert form; quantised, $Z$ generates the twistor attachment's conformal energy $\tfrac12(N_{a} + N_{b}) + 1 = |s| + 1 + n$ with the stored multiplicities, future translations, null special conformal transformations and the de Sitter horizon generator are nonnegative with it, and the attachment's Fock vacuum satisfies the positivity and invariance clause of (N$_{+}$) on Minkowski local horizons (and de Sitter for conformal fields) with no thermal input; on those horizons (N$_{+}$) becomes the choice of that vacuum, a restatement and not a derivation of the premise; the eternal hole is not reached.
\emph{Source:} Transport Lab claim synthesis.attachment-conformal-energy-is-generated-by-z, version 1.
\end{fact}

\begin{fact}[On the de Sitter gravitational field $\hb = \sigma$ on flat spacetime the uniform affine translation along a cosmological horizon's generators is the background conformal null rotation generated by $\tfrac12 nI = P_{n} + \tfrac{\Lambda}{12}\,n\nO$, which nonnegative conformal energy makes nonnegative; for a conformally invariant field in its conformal vacuum this gives, through an argued identification of the horizon modes, the positivity and invariance clause of (N$_{+}$) case (b); the Schwarzschild field, whose field strength is pure Weyl, admits no such gauge]\label{c:synthesis.de-sitter-horizon-translation-is-conformal-null-rotation}
On the de Sitter gravitational field $\hb = \sigma$ on flat spacetime a cosmological horizon is a background null plane whose lines are free-fall null worldlines with $\sigma$ constant along them, and the uniform shift of the physical affine parameter $s/\sigma$ is the background conformal null rotation generated by $\tfrac12 nI = P_{n} + \tfrac{\Lambda}{12}n\nO$, which the stored dilation and boost limits place in the cone of nonnegative generators; so nonnegative conformal energy gives the positivity half of (N$_{+}$) on every de Sitter cosmological horizon for a conformally invariant field in its conformal vacuum (argued identification of the horizon modes), while the Schwarzschild field, whose field strength is pure nonzero Weyl, admits no gauge in which the argument runs.
\emph{Source:} Transport Lab claim synthesis.de-sitter-horizon-translation-is-conformal-null-rotation, version 4.
\end{fact}

\begin{fact}[Under hypothesis (F), that the descended universe's late de Sitter stage is the frozen point's own field, a deposit correlated by the frozen point's stabiliser has range beyond the stored requirement by over $4 \times 10^{5}$ but reads every observed mode as inside the frozen point's horizon, so its spectrum is the vacuum's, $n_{s} - 1 = 2$]\label{c:synthesis.frozen-point-geometry-reads-descent-as-vacuum}
Under hypothesis (F) (the late de Sitter stage is the frozen point's own field), a deposit correlated by the frozen point's stabiliser invariant $Z$ reaches $|x_{d}| = \eta_{\infty} - \eta_{d}$, $a_{0}H_{0}|x_{d}| \in (4.32, 4.34)$, which exceeds the required range ($a_{0}H_{0}R_{\mathrm{req}} \approx 9.5 \times 10^{-6}$, one length for every $T_{d}$) by over $4 \times 10^{5}$; but with $k_{*}|x_{d}| \in (960, 968)$ every observed mode is inside the frozen point's horizon at the descent, so the invariant (Euclidean) state of every unitary-series field, and the massless Bunch-Davies state, give the vacuum's $n_{s} - 1 = 2$ up to $O(u^{-2})$, more than $480$ standard deviations from the measured value: range is not the obstruction, horizon exit is.
\emph{Source:} Transport Lab claim synthesis.frozen-point-geometry-reads-descent-as-vacuum, version 1.
\end{fact}

\begin{fact}[Imported: a continuous positive-definite function is the Fourier transform of a finite positive measure (Bochner); a derivative of even order $2k$ of that transform at the origin gives a finite moment of order $2k$ (the moment lemma, Lukacs); a smooth even function of one variable is a smooth function of its square (Whitney 1943)]\label{c:synthesis.imported-bochner-theorem-moment-lemma-and-whitney-even-functions}
Bochner: continuous positive-definite functions on $\R$ are the Fourier transforms of finite positive measures; the moment lemma (Lukacs): a derivative of order $2k$ of the transform at $0$ gives a finite moment of order $2k$; Whitney (1943): a smooth even function is $g(x^{2})$ with $g$ smooth, a smooth odd one $x\,g(x^{2})$.
\emph{Source:} S. Bochner, Monotone Funktionen, Stieltjessche Integrale und harmonische Analyse, Math. Ann. 108 (1933), 378--410.; E. Lukacs, Characteristic Functions, 2nd ed., Griffin, London (1970)., ch. 2; H. Whitney, Differentiable even functions, Duke Math. J. 10 (1943), 159--160..
\end{fact}

\begin{fact}[Imported: a four-dimensional gravitational field with vanishing Weyl curvature is locally a conformal rescaling of the flat field (Weyl, Schouten), and in four dimensions the conformally coupled massless scalar (weight 1), the massless Dirac field (weight 3/2) and the Maxwell and Yang-Mills fields (weight 0) are conformally covariant]\label{c:synthesis.imported-conformal-covariance-of-massless-fields-and-local-conformal-flatness}
A four-dimensional gravitational field with vanishing Weyl curvature is locally $\hb(b) = b/\varphi$, the flat field rescaled (Weyl, Schouten); under $g \mapsto \varphi^{2}g$ the conformally coupled massless scalar rescales by $\varphi^{-1}$, the massless Dirac field by $\varphi^{-3/2}$, and the Maxwell and Yang-Mills equations are unchanged, with the field algebras identified.
\emph{Source:} J. A. Schouten, Ueber die konforme Abbildung n-dimensionaler Mannigfaltigkeiten mit quadratischer Massbestimmung auf eine Mannigfaltigkeit mit euklidischer Massbestimmung, Math. Z. 11 (1921), 58--88.; H. Weyl, Reine Infinitesimalgeometrie, Math. Z. 2 (1918), 384--411.; R. M. Wald, General Relativity, University of Chicago Press (1984)., Appendix D; R. Penrose and W. Rindler, Spinors and Space-Time, Vol. 1, Cambridge University Press (1984)., Sec. 5.7.
\end{fact}

\begin{fact}[Imported: for a massive scalar field on de Sitter the de Sitter-invariant states form the alpha-family, of which only the Euclidean state is Hadamard; the others are singular at antipodal points (Allen 1985)]\label{c:synthesis.imported-de-sitter-invariant-hadamard-state-unique}
For $m^{2} > 0$ the de Sitter-invariant states are the $\alpha$-vacua, Bogoliubov transforms of the Euclidean state; only the Euclidean state is Hadamard, the others being singular at antipodal points (Allen 1985).
\emph{Source:} B. Allen, Vacuum states in de Sitter space, Phys. Rev. D 32 (1985), 3136--3149.; R. Brunetti, K. Fredenhagen and S. Hollands, A remark on alpha vacua for quantum field theories on de Sitter space, J. High Energy Phys. 2005, no. 05, 063; arXiv:hep-th/0503022..
\end{fact}

\begin{fact}[Imported: on every gravitational field with a globally hyperbolic derived metric the free Dirac field exists as field operators with a causal anticommutator, and its CAR algebras of local observables satisfy a generalised form of the Haag-Kastler axioms; the causal propagator reproduces every solution from its Cauchy data, and the anticommutator pairing is the conserved Cauchy inner product of the solutions the test spinors define (Dimock 1982)]\label{c:synthesis.imported-dimock-dirac-car-algebra-on-a-gravitational-field}
On every gravitational field with a globally hyperbolic derived metric the free Dirac field exists, with a causal anticommutator and CAR algebras of local observables satisfying a generalised form of the Haag-Kastler axioms; the causal propagator reproduces every solution from its Cauchy data and the anticommutator pairing is the conserved Cauchy inner product of the solutions; no state is selected (Dimock 1982).
\emph{Source:} J. Dimock, Dirac quantum fields on a manifold, Trans. Amer. Math. Soc. 269 (1982), 133--147., Sec. II.A (S = S+ - S-, after Thm. 2.1); Sec. II.B (Prop. 2.2, Thm. 2.3, Prop. 2.4(b)-(d)); Sec. III.A (the inner product on Cauchy data); Sec. III.B and Thm. 3.1 (the anticommutator).
\end{fact}

\begin{fact}[Imported: for the Dirac field the Hadamard form is equivalent to the wave front set condition and propagates from a neighbourhood of one Cauchy surface to the whole field (Kratzert 2000; Hollands 2001; Sahlmann and Verch 2001), on Robertson-Walker fields the Dirac states adiabatic of infinite order are Hadamard (Hollands 2001), and quasi-free Hadamard states have as normal-coordinate scaling limit the flat massless Dirac two-point function (Sahlmann and Verch 2001, Prop. 5.6, posed for Majorana spinors)]\label{c:synthesis.imported-dirac-hadamard-form-propagates-and-adiabatic-states-are-hadamard}
For the free Dirac field on a globally hyperbolic gravitational field, the Hadamard form is equivalent to the wave front set condition on both orderings and propagates from a neighbourhood of one Cauchy surface to the whole field (Kratzert 2000; Hollands 2001; Sahlmann and Verch 2001); on Robertson-Walker fields homogeneous Dirac states adiabatic of infinite order are Hadamard (Hollands 2001; the converse is not imported); quasi-free Hadamard states scale in normal coordinates to the flat massless Dirac two-point function (Sahlmann and Verch 2001, Prop. 5.6, posed for Majorana spinors; Sanders 2010 for the Dirac equivalence).
\emph{Source:} K. Kratzert, Singularity structure of the two point function of the free Dirac field on a globally hyperbolic spacetime, Ann. Phys. (Leipzig) 9 (2000), 475--498; arXiv:math-ph/0003015.; S. Hollands, The Hadamard condition for Dirac fields and adiabatic states on Robertson-Walker spacetimes, Commun. Math. Phys. 216 (2001), 635--661; arXiv:gr-qc/9906076.; H. Sahlmann and R. Verch, Microlocal spectrum condition and Hadamard form for vector-valued quantum fields in curved spacetime, Rev. Math. Phys. 13 (2001), 1203--1246; arXiv:math-ph/0008029., Sec. 5.2; Thm. 5.8; Def. 5.3; Prop. 5.6; Lemma 5.7; K. Sanders, The locally covariant Dirac field, Rev. Math. Phys. 22 (2010), 381--430; arXiv:0911.1304..
\end{fact}

\begin{fact}[Imported: on an ultrastatic slab with compact spatial sections the pure quasi-free state of the free Dirac field extracted from the fermionic projector, the positive part of the weighted spacetime spin product, generically fails to be Hadamard, while the softened states built with a smooth compactly supported time weight are Hadamard (Fewster and Lang 2015)]\label{c:synthesis.imported-fewster-lang-spin-product-states-of-the-dirac-field}
On an ultrastatic slab with compact sections the Dirac state from the positive part of the weighted spin product (the fermionic projector state) is pure and quasi-free, generically not Hadamard with a sharp time weight, and Hadamard with a smooth compactly supported weight (Fewster-Lang 2015, Thms 3.4, 3.5).
\emph{Source:} C. J. Fewster and B. Lang, Pure quasifree states of the Dirac field from the fermionic projector, Class. Quantum Grav. 32 (2015), 095001; arXiv:1408.1645., Sec. 3: the state from the positive spectral projection of the weighted spacetime spin product; Thm. 3.4 (the unsoftened state is not Hadamard for generic lifetimes); Thm. 3.5 (softened states are Hadamard); F. Finster and M. Reintjes, A non-perturbative construction of the fermionic projector on globally hyperbolic manifolds I: space-times of finite lifetime, Adv. Theor. Math. Phys. 19 (2015), 761--803; arXiv:1301.5420..
\end{fact}

\begin{fact}[Imported: on spatially flat Robertson-Walker fields a homogeneous scalar state whose modes grow as $k^{-\nu}$ with $\nu \ge \tfrac32$ at small momentum has an infrared-divergent two-point function, so the massless minimally coupled Bunch-Davies state is not a state on slice fields of nonzero mean (Ford and Parker 1977)]\label{c:synthesis.imported-ford-parker-infrared-divergence}
On spatially flat Robertson-Walker fields, if a homogeneous scalar state's modes grow as $k^{-\nu}$, $\nu \ge \tfrac32$, at small momentum, it has an infrared-divergent two-point function, as the massless minimally coupled Bunch-Davies state does (Ford and Parker 1977): slice fields of nonzero mean have infinite variance, those of zero mean finite variance for $\nu < \tfrac52$.
\emph{Source:} L. H. Ford and L. Parker, Infrared divergences in a class of Robertson-Walker universes, Phys. Rev. D 16 (1977), 245--250.; A. Higuchi and N. Rendell, Infrared divergences for free quantum fields in cosmological spacetimes, arXiv:1711.03964., Sec. 3.
\end{fact}

\begin{fact}[Imported: the massless two-point function $D_{+}$, real and even at spacelike separation; Pauli's free-field spin-statistics argument (1940); and the algebraic spin-statistics theorem of Guido and Longo (1995): for a modular covariant net of field algebras with the Reeh-Schlieder property and normal commutation relations the statistics grading is $\Gamma = U(2\pi)$]\label{c:synthesis.imported-free-field-spin-statistics-and-guido-longo-theorem}
$D_{+}(\xi) = -1/(4\pi^{2}(Q(\xi) - i0\,\xi_{0}))$, real and even at spacelike $\xi$; Pauli's free-field spin-statistics argument; and Guido-Longo: modular covariance, Reeh-Schlieder and normal commutation relations give $\Gamma = U(2\pi)$ with positive energy derived.
\emph{Source:} W. Pauli, The connection between spin and statistics, Phys. Rev. 58 (1940), 716--722.; D. Guido and R. Longo, An algebraic spin and statistics theorem, Commun. Math. Phys. 172 (1995), 517, arXiv:funct-an/9406005, Theorem 2.11 and Proposition 2.4.; A. Shalabi, M. H. Zambianco and T. R. Perche, The effect of curvature on local observables in quantum field theory, arXiv:2412.12294, equation (11).; R. F. Streater and A. S. Wightman, PCT, Spin and Statistics, and All That, Benjamin (1964)..
\end{fact}

\begin{fact}[Imported: the dictionary between Gauge Theory Gravity and curved-spacetime language]\label{c:synthesis.imported-gtg-dictionary}
On the invertible region, tangent vectors are $\hb^{-1}(a)$, the Levi-Civita derivative is $D_a$, geodesics have $D$-parallel velocity, and the Riemann, Ricci, Schouten and Weyl tensors and Killing fields are read from the field strength $\mathcal{R}(B)$ and $D$.
\emph{Source:} A. Lasenby, C. Doran and S. Gull, Gravity, gauge theories and geometric algebra, Phil. Trans. R. Soc. Lond. A 356 (1998), 487--582., \S\S 3.5, 4.3--4.4, 5.1--5.2, App.~C; C. Doran and A. Lasenby, Geometric Algebra for Physicists, Cambridge University Press (2003)., \S\S~13.5--13.6, 14.1; A. N. Lasenby, Geometric algebra, gravity and gravitational waves, Adv. Appl. Clifford Algebras (2019)., \S\S~4--5, eq.~(8.5) (arXiv:1912.05960); R. M. Wald, General Relativity, University of Chicago Press (1984)., \S~12.5.
\end{fact}

\begin{fact}[Imported: the Hadamard parametrix of a normally hyperbolic operator on a vector bundle has leading coefficient the square root of the Van Vleck determinant times the parallel transport of the operator's connection (Friedlander 1975; Bar, Ginoux and Pfaffle 2007, Ch. 2); for the squared Dirac operator that connection is the spin connection (Lichnerowicz), so U = Delta\^{}(1/2) R; the Hadamard form does not depend on the time function of its i0 prescription (Kay and Wald 1991; Sahlmann and Verch 2001, Sec. 5.1); and in four dimensions the retarded and advanced fundamental solutions on a causal domain are (1/4 pi) theta(+-(t - t')) [U delta(sigma) + V theta(sigma)] with V smooth (Friedlander 1975; Bar, Ginoux and Pfaffle 2007, Ch. 2)]\label{c:synthesis.imported-hadamard-parametrix-of-the-squared-dirac-operator-leads-with-van-vleck-root-times-spin-transport}
For a normally hyperbolic operator on a vector bundle the leading Hadamard coefficient is $\Delta^{1/2}$ times the parallel transport of the operator's connection (Friedlander 1975; Bar, Ginoux and Pfaffle 2007); for the squared Dirac operator that is the spin connection, so $U = \Delta^{1/2}R$ with $\dot R = -\tfrac12\Omega(\dot\gamma)R$; the Hadamard form does not depend on the time function of its $i0$ prescription (Kay and Wald 1991; Sahlmann and Verch 2001, Sec. 5.1); in four dimensions the retarded and advanced fundamental solutions on a causal domain are $\tfrac{1}{4\pi}\theta(\pm(t - t'))[U\delta(\sigma) + V\theta(\sigma)]$, $V$ smooth.
\emph{Source:} F. G. Friedlander, The Wave Equation on a Curved Space-Time, Cambridge University Press, Cambridge, 1975., the fundamental solutions of normally hyperbolic operators and their Hadamard form; cited from memory, not read; C. Bar, N. Ginoux and F. Pfaffle, Wave Equations on Lorentzian Manifolds and Quantization, ESI Lectures in Mathematics and Physics, European Mathematical Society, 2007; arXiv:0806.1036., Ch. 2 (Def. 2.2.1, Prop. 2.3.1; the formal and local fundamental solutions as sums of Hadamard coefficients times Riesz distributions); Sec. 1.3 (the function mu\_x); E. Poisson, The motion of point particles in curved spacetime, Living Rev. Relativity 7 (2004), 6; arXiv:gr-qc/0306052., Sec. 4.3.2, Eqs. (266), (268), (269), (273) (scalar Green functions); Sec. 4.4.2 (electromagnetic Green functions, U = parallel propagator times Delta\^{}(1/2)); H. Sahlmann and R. Verch, Microlocal spectrum condition and Hadamard form for vector-valued quantum fields in curved spacetime, Rev. Math. Phys. 13 (2001), 1203--1246; arXiv:math-ph/0008029., Sec. 3.5; Sec. 5.1; B. S. Kay and R. M. Wald, Theorems on the uniqueness and thermal properties of stationary, nonsingular, quasifree states on spacetimes with a bifurcate Killing horizon, Phys. Rep. 207 (1991), 49--136.; M. J. Radzikowski, Micro-local approach to the Hadamard condition in quantum field theory on curved space-time, Commun. Math. Phys. 179 (1996), 529--553.; L. Ronge, Formulas for Hadamard coefficients in terms of Green's operators, arXiv:2604.25763 (2026)., Def./Prop. 2.9, 2.10; Sec. 2.1; A. G. M. Lewis, Hadamard renormalization of a 2-dimensional Dirac field, Phys. Rev. D 101 (2020), 125019; arXiv:1910.05663., Eqs. (14), (26), (52).
\end{fact}

\begin{fact}[Imported: on an arbitrary four-dimensional spacetime, near the vertex of the future light cone of a point (inside a geodesically convex, globally hyperbolic neighbourhood), the trace of a Dirac field on the cone, its two transverse spinor components, is an L2 isometry from data on a spacelike slice (the flux identity of the conserved current), the characteristic Cauchy problem on the cone is uniquely solvable for L2 data, and the trace is an isomorphism of H1 onto the space of cone data (Haefner and Nicolas 2011, Theorems 1 and 2)]\label{c:synthesis.imported-haefner-nicolas-dirac-characteristic-problem-on-light-cones-near-the-vertex}
For the Dirac field of any mass on an arbitrary four-dimensional spacetime, near the vertex of the future light cone of a point (in a geodesically convex, globally hyperbolic neighbourhood, below a level $t = T$), the trace on the cone, two transverse spinor components, is an $L^{2}$ isometry from data on a spacelike slice by the flux identity of the conserved current, the characteristic problem is uniquely solvable for $L^{2}$ data, so the trace is unitary, and it is an isomorphism of $H^{1}$ onto the cone data (Häfner and Nicolas 2011, Theorems 1 and 2).
\emph{Source:} D. Häfner and J.-P. Nicolas, The characteristic Cauchy problem for Dirac fields on curved backgrounds, J. Hyperbolic Differ. Equ. 8 (2011), no. 3, 437--483; arXiv:0903.0515., Abstract; Sec. 3 (Def. 3.6, the domain D\_T); Sec. 4.1 (Def. 4.2, Eq. (14), Thm. 1); Sec. 4.4 (Def. 4.11, Thm. 2).
\end{fact}

\begin{fact}[Imported: on a spacetime with a bifurcate Killing horizon there is at most one stationary quasi-free Hadamard state, KMS at the Hawking temperature on a wedge; none exists on extended Schwarzschild-de Sitter, and on Kerr none is argued (Kay and Wald 1991, with Kay and Lupo 2016)]\label{c:synthesis.imported-kay-wald-uniqueness-and-no-go}
On a spacetime with a bifurcate Killing horizon there is at most one stationary quasi-free Hadamard state, KMS at $\kappa/2\pi$ on a wedge (Kay and Wald 1991); none exists on extended Schwarzschild-de Sitter, and on Kerr none is argued for the scalar field; a gap in the proofs was filled by Kay and Lupo (2016, App. B).
\emph{Source:} B. S. Kay and R. M. Wald, Theorems on the uniqueness and thermal properties of stationary, nonsingular, quasifree states on spacetimes with a bifurcate Killing horizon, Phys. Rep. 207 (1991), 49--136.; B. S. Kay and U. Lupo, Non-existence of isometry-invariant Hadamard states for a Kruskal black hole in a box and for massless fields on 1+1 Minkowski spacetime with a uniformly accelerating mirror, Class. Quantum Grav. 33 (2016), 215001; arXiv:1502.06582., App. B; footnotes 2 and 4; B. S. Kay, Quantum field theory in curved spacetime (2nd edition), in Encyclopedia of Mathematical Physics, 2nd edition (2024); arXiv:2308.14517..
\end{fact}

\begin{fact}[Imported: quasi-free Hadamard states are locally quasi-equivalent, for the Klein-Gordon field (Verch 1994) and for the Dirac field (D'Antoni and Hollands 2006)]\label{c:synthesis.imported-local-quasi-equivalence-of-hadamard-states}
Any two quasi-free Hadamard states are locally quasi-equivalent on the algebra of every open relatively compact region, for the Klein-Gordon field (Verch 1994) and for the Dirac field, there with local unitary equivalence and type III$_{1}$ local algebras (D'Antoni and Hollands 2006).
\emph{Source:} R. Verch, Local definiteness, primarity and quasiequivalence of quasifree Hadamard quantum states in curved spacetime, Commun. Math. Phys. 160 (1994), 507--536.; C. D'Antoni and S. Hollands, Nuclearity, local quasiequivalence and split property for Dirac quantum fields in curved spacetime, Commun. Math. Phys. 261 (2006), 133--159; arXiv:math-ph/0106028., Thms. III.1, III.2.
\end{fact}

\begin{fact}[Imported: light-front quantisation of free fields: on a null plane the free fields of different masses are unitarily equivalent and generate an irreducible algebra (Leutwyler, Klauder and Streit 1970), and the restriction of a free field to a null surface carries no mass except in zero modes, the mass entering the Dirac field only through its constrained half (Wall 2012)]\label{c:synthesis.imported-null-plane-restriction-of-free-fields-is-mass-independent}
On a null plane free fields of different masses are unitarily equivalent and generate an irreducible algebra (Leutwyler, Klauder and Streit 1970); restricted to a null surface a free field's algebra and null energy carry no mass except in zero modes, the Dirac field's mass entering only its constrained half (Wall 2012).
\emph{Source:} H. Leutwyler, J. R. Klauder and L. Streit, Quantum field theory on lightlike slabs, Nuovo Cimento A 66 (1970), 536--554.; A. C. Wall, A proof of the generalized second law for rapidly changing fields and arbitrary horizon slices, Phys. Rev. D 85 (2012), 104049; arXiv:1105.3445., Secs. 3.2, 4.1.
\end{fact}

\begin{fact}[Imported: quasi-equivalence of quasi-free CAR states and the Powers-Stormer inequality (Powers-Stormer), and disjointness of KMS states at different temperatures (Takesaki)]\label{c:synthesis.imported-powers-stormer-takesaki-disjointness}
Gauge-invariant quasi-free CAR states $\omega_{A}$, $\omega_{B}$ are factor states, quasi-equivalent iff $A^{1/2} - B^{1/2}$ and $(1 - A)^{1/2} - (1 - B)^{1/2}$ are Hilbert-Schmidt and otherwise disjoint, and $\|A^{1/2} - B^{1/2}\|_{2}^{2} \le \|A - B\|_{1}$ for bounded positive $A$, $B$ (Powers-Størmer); KMS states at different temperatures are disjoint when one induces a type III algebra (Takesaki).
\emph{Source:} R. T. Powers and E. Stormer, Free states of the canonical anticommutation relations, Commun. Math. Phys. 16 (1970), 1--33.; M. Takesaki, Disjointness of the KMS-states of different temperatures, Commun. Math. Phys. 17 (1970), 33--41..
\end{fact}

\begin{fact}[Imported: the Hadamard form of a two-point function in a neighbourhood of one Cauchy surface propagates to the whole spacetime (Fulling, Sweeny and Wald 1978)]\label{c:synthesis.imported-propagation-of-hadamard-form}
For the free scalar field, a two-point function of Hadamard form in a neighbourhood of one Cauchy surface is of Hadamard form everywhere: Cauchy evolution preserves the Hadamard form (Fulling, Sweeny and Wald 1978).
\emph{Source:} S. A. Fulling, M. Sweeny and R. M. Wald, Singularity structure of the two-point function in quantum field theory in curved spacetime, Commun. Math. Phys. 63 (1978), 257--264..
\end{fact}

\begin{fact}[On a horizon's null plane the canonical part of a spacetime-algebra spinor is the image of Lambda(m), two complex channels, one per chirality, with helicity about the generators equal to minus the chirality; so, given one (Q) version 3 array per complex channel, one massless Dirac species supplies two arrays per light ray (ln 2 at the Hawking temperature) and one Weyl field one]\label{c:synthesis.light-ray-good-part-of-a-spinor-is-two-chirality-arrays}
In $\Cl(1,3)$ the flux of the Dirac current through a horizon's null plane $N_{k}$ is $2|\Lambda(m)\psi|^{2}$, so the canonical part there is the image of $\Lambda(m) = m\gamma_{0}/2$, real dimension $4$, of boost weight $\tfrac12$ in the carrier's convention, split by chirality into two complex channels with helicity about $k$ minus the chirality; given one (Q) version 3 array per channel, one massless Dirac species gives two arrays per light ray, $\ln 2$ at the Hawking temperature, a Weyl field one, and the per-light-ray total over fields counts two arrays per registered unit.
\emph{Source:} Transport Lab claim synthesis.light-ray-good-part-of-a-spinor-is-two-chirality-arrays, version 2.
\end{fact}

\begin{fact}[On a horizon generator the weight-$\tfrac12$ scaling preserves the Hilbert-Schmidt norm of a kernel shrunk onto a point, so a gauge-invariant quasi-free state locally normal to the null vacuum at a point is locally the null vacuum there in every regular parameter, and a stationary gauge-invariant quasi-free state locally normal to it at the edge is the null vacuum of (N$_{+}$), positivity included; every other stationary gauge-invariant quasi-free state, the half-thermal products at every period included, is locally disjoint from it at the edge]\label{c:synthesis.local-normality-at-the-edge-gives-the-null-ground-state}
On a horizon generator the weight-$\tfrac12$ scaling preserves the Hilbert-Schmidt norm of a kernel shrunk onto a point, so for gauge-invariant quasi-free states local normality to the null vacuum at a point (via Powers-Stormer, a square-integrable difference) gives the local form of (N$_{+}$) there in any regular parameter, and a stationary such state locally normal to it at the edge is the null vacuum of (N$_{+}$), positivity included; every other stationary gauge-invariant quasi-free state differs by a degree $-1$ kernel that is not square integrable at the edge, so the half-thermal products, $c = 1$ included, are locally disjoint from the vacuum there.
\emph{Source:} Transport Lab claim synthesis.local-normality-at-the-edge-gives-the-null-ground-state, version 2.
\end{fact}

\begin{fact}[A ground state for the affine null translations of a horizon's generators restricts on either side of the edge to the continuum I-state at rapidity period $2\pi$, so the continuum restatements of (R) for the eternal hole, (R$_{\Lambda}$), (R$_{\mathrm{loc}}$) and the joint KMS part of (Q) follow from one premise with no temperature in it, posable only on the continuum]\label{c:synthesis.positive-null-energy-gives-r-premises}
A gauge-invariant quasi-free ground state for the affine null translations of a horizon's generators (premise (N$_{+}$), no temperature) restricts on either side of the edge to the continuum I-state at rapidity period exactly $2\pi$ for the boost future-directed there. So the continuum restatements of (R) for the eternal hole (a Hartle-Hawking-type regularity, collapse not covered), of (R$_{\Lambda}$) and of (R$_{\mathrm{loc}}$) in Minkowski space, with the joint KMS part of (Q), follow from one premise. The stored point-spectrum wordings are its discretisations, and the import is relocated to regularity across the bifurcation surface.
\emph{Source:} Transport Lab claim synthesis.positive-null-energy-gives-r-premises, version 4.
\end{fact}

\begin{fact}[A quasi-free deposit on the fermionic boundary-mode algebra is its vacuum plus shot noise, whatever its correlation range or partner placement; for the I-state, or any quasi-free deposit with at least one thermal species' number of quanta at thermal energies, the shot term at the pivot is below $10^{-77}$]\label{c:synthesis.quasi-free-deposit-is-vacuum-plus-shot-noise}
For every quasi-free state of a CAR algebra and every bilinear observable, $\mathrm{sd}_{\omega}(O) \le \mathrm{sd}_{\omega_{0}}(O) + \Tr(F^{2}X)^{1/2}$ with $X = (\Gamma - \Gamma_{0})^{2} + \Gamma(1 - \Gamma)$, $\Tr X = 2\omega(N_{0})$: if a deposit is quasi-free it is the vacuum plus shot noise whatever its correlation range or partner placement, and the descent's parity average adds a region-independent constant; for the I-state, or any occupation profile with at least one thermal species' number of quanta and energy spread $c_{k} \le 1.31$, the shot term at the pivot is below $10^{-77}$ for every reheating temperature, far below $A_{s} = 2.1 \times 10^{-9}$, while a sparse or energetically broad deposit is bounded only by $(3c_{k}k^{3}/2\pi^{2}\bar n)^{1/2}$.
\emph{Source:} Transport Lab claim synthesis.quasi-free-deposit-is-vacuum-plus-shot-noise, version 3.
\end{fact}

\begin{fact}[The light-ray arrays of the slot-module matter are its Weyl components, sixteen per generation for every CPT halving and every mass term, so a generation registers eight units of ln 2; the particle half is chiral in every halving, so the counted species are not gauge-covariant Dirac species, and with a nonzero Majorana term not Dirac species by mass either]\label{c:synthesis.slot-module-light-ray-arrays-are-weyl-components}
On a horizon's null plane mass terms reach only the constraint, so the arrays of a light ray are the Weyl components; every CPT halving of the slot module has sixteen per generation (eight registered units, for either neutral singlet), and its particle half is chiral in every halving, so the series' matter cannot be grouped into gauge-covariant Dirac species and, with a nonzero Majorana term, is seven Dirac species and two Majorana neutrals per generation.
\emph{Source:} Transport Lab claim synthesis.slot-module-light-ray-arrays-are-weyl-components, version 1.
\end{fact}

\begin{fact}[On a horizon generator the boost is the anticommutator-preserving scaling about the edge, so a stationary gauge-invariant quasi-free state is its own scaling limit there and is the null vacuum of (N$_{+}$) exactly when it is locally the null vacuum at the edge; stationary states thermal on each half at any period, with or without cross-edge correlation, are locally the null vacuum at every other point, so the content of (N$_{+}$) beyond stationarity sits at the bifurcation surface alone]\label{c:synthesis.stationary-horizon-state-is-its-own-edge-scaling-limit}
On every horizon in the scope of (N$_{+}$) the boost is the anticommutator-preserving scaling of a generator about its edge, so a stationary quasi-free state is its own short-distance limit there and satisfies (N$_{+}$) exactly when it is locally the null vacuum at the edge; the stationary states thermal on each half at any period (and at the right period with no cross-edge correlation) are locally the null vacuum at every other point, and zero null energy does not select the vacuum either, so within the stationary class the content of (N$_{+}$) is one temperature-free condition at the bifurcation surface, which on the eternal hole lies off the Painlev\'e-Gullstrand chart.
\emph{Source:} Transport Lab claim synthesis.stationary-horizon-state-is-its-own-edge-scaling-limit, version 2.
\end{fact}

\begin{fact}[A sudden descent into the frozen phase with an instantaneous-vacuum start leaves a feature, not a red tilt]\label{c:synthesis.sudden-descent-feature-not-tilt}
If the descent is a single event at $\eta_d$ into the frozen phase and a massless field starts there in the instantaneous vacuum, its spectrum is $(H/2\pi)^{2}F(k|\eta_d|)$ with $F \approx \tfrac49 x^{2}$ outside and $F \to 1$ with $x^{-2}$ ringing inside, and no seven-e-fold window carries the measured red index (exact for every window); such a sudden start does not supply the tilt.
\emph{Source:} Transport Lab claim synthesis.sudden-descent-feature-not-tilt, version 3.
\end{fact}

\begin{fact}[On a globally hyperbolic Weyl-flat region of a gravitational field every quasi-free Hadamard state of the massless Dirac field has the flat massless vacuum as its scaling limit and, with the comoving transverse labels, the generator microlocal condition in the physical affine parameter, so the local form of (N$_{+}$), at the leading order at which local Rindler horizons are defined, holds there for the massless field without the scaling-limit premise; on a Robertson-Walker field the conformal vacuum differs on every null generator from the null vacuum of the physical affine parameter by the Schwarzian kernel of $S = 2a''/a - 4(a'/a)^{2}$ (conformal time), which vanishes identically exactly on de Sitter and flat fields]\label{c:synthesis.weyl-flat-fields-need-no-scaling-limit-premise-and-robertson-walker-generators-see-the-schwarzian}
On a globally hyperbolic Weyl-flat region every quasi-free Hadamard state of the massless Dirac field has the flat massless vacuum as its scaling limit and, with comoving transverse labels, the generator microlocal condition in the physical affine parameter, so for the massless field the local form of (N$_{+}$) (at the leading order defining local Rindler horizons) needs no scaling-limit premise there; on a Robertson-Walker field the conformal vacuum differs on null generators from the physical null vacuum by a Schwarzian kernel with $S = 2a''/a - 4(a'/a)^{2}$ (leading coefficient $-S/12a^{4}$ in the physical parameter), zero exactly on de Sitter and flat fields (radiation $-4/t^{2}$).
\emph{Source:} Transport Lab claim synthesis.weyl-flat-fields-need-no-scaling-limit-premise-and-robertson-walker-generators-see-the-schwarzian, version 4.
\end{fact}

\section*{Conclusion}
The Hawking temperature is not built into the algebra: it enters through the geometric period, and under (P) and (R) it becomes the temperature of every boundary mode. The radiation side also rests on the interface identification (I), with the further premise (Q) for the boson-like part of the spectrum.

The thermal premises are now gathered into (N$_{+}$), a premise about the state with no temperature in it. Its local form is settled for the Dirac field of any mass: every gauge-invariant quasi-free Hadamard state meets it at every point of every local Rindler horizon. That the series' states are Hadamard is the adopted premise (HD), and a thermal deposit on a smooth descent of radiation type would meet it.

Two questions remain open. The first is which state: the Hadamard condition leaves the deposit's spectrum free on every bounded band, so nothing it implies fixes the occupation profile. Selecting it by the Dirac field's own spin product fails for massless fields.

The second is the ground-state clause of (N$_{+}$), which no Hadamard condition gives and which needs a global datum such as stationarity at the edge. Its equivalent form, thermality of the future half beyond every point, would follow from geometric modular action, which reaches Minkowski null planes and de Sitter horizons but not the eternal hole.

\end{document}